%g16.tex

\documentclass{amsart}
\usepackage{amssymb,amsfonts,latexsym,amscd}

\usepackage[all]{xy}
\usepackage{verbatim}
\usepackage{hyperref}
\usepackage{mathrsfs}
\usepackage{mathtools}
\usepackage{amssymb}
\usepackage{stmaryrd}
\usepackage{tikz-cd}
\usepackage{xcolor}

\newcommand{\gs}{\textcolor{blue}}

\newtheorem{theorem}{Theorem}[section]
\newtheorem{lemma}[theorem]{Lemma}
\newtheorem{proposition}[theorem]{Proposition}
\newtheorem{corollary}[theorem]{Corollary}
\theoremstyle{definition}
\newtheorem{definition}[theorem]{Definition}

\newtheorem{example}[theorem]{Example}
\newtheorem{question}[theorem]{Question}
\newtheorem{conjecture}[theorem]{Conjecture}
\newtheorem{remark}[theorem]{Remark}
\newtheorem{claim}[theorem]{Claim}
\newtheorem{note}[theorem]{Note}
\newtheorem{thm}{Theorem}[subsection]
\newtheorem{lem}[thm]{Lemma}
\newtheorem{prop}[thm]{Proposition}
\newtheorem{cor}[thm]{Corollary}
\newtheorem{rem}[thm]{Remark}
\newtheorem{exam}[thm]{Example}
\newcommand{\id}{{\rm id}}
\newcommand{\End}{{\rm End}}
\newcommand{\Fun}{{\rm Fun}}
\newcommand{\Irr}{\text{\rm Irr}}
\newcommand{\Pro}{\text{Pro}}
\newcommand{\Hom}{{\rm Hom}}
\newcommand{\Ad}{{\rm Ad}}
\newcommand{\Id}{{\rm Id}}
\newcommand{\Comod}{{\rm Comod}}
\newcommand{\Aut}{{\rm Aut}}
\newcommand{\Rep}{{\rm Rep}}
\newcommand{\Corep}{{\rm Corep}}
\newcommand{\Vect}{{\rm Vec}}
\renewcommand{\O}{\mathscr{O}}
\newcommand{\C}{\mathscr{C}}
\newcommand{\M}{\mathscr{M}}
\newcommand{\V}{\mathscr{V}}
\newcommand{\la}{\langle\,}
\newcommand{\ra}{\,\rangle}
\newcommand{\ot}{\otimes}
\newcommand{\qexp}{{\text{qexp}}}
\newcommand{\ben}{\begin{enumerate}}
\newcommand{\een}{\end{enumerate}}
\newcommand{\ad}{{\text{ad\,}}}
\newcommand{\Rad}{{\text{Rad}}}
\newcommand{\GL}{{\text{GL}}}
\newcommand{\Lie}{{\text{Lie}}}
\newcommand{\D}{{\mathcal D}}
\newcommand{\cA}{{\mathcal A}}
\newcommand{\E}{\mathcal{E}}
\newcommand{\Z}{{\mathcal Z}}

\hyphenation{se-mi-simple co-se-mi-simple}

\numberwithin{equation}{section}

\begin{document}

\title[Hopf $2$-cocycles for certain affine algebraic reductive groups] {Hopf $2$-cocycles for certain affine algebraic reductive groups}

\author{Shlomo Gelaki}
\address{Department of Mathematics, Iowa State University, Ames, IA 50011, USA} 
\email{gelaki@iastate.edu}

\date{\today}

\keywords{affine algebraic reductive groups; semisimple tensor categories;
classical fiber functors; Hopf $2$-cocycles}

\begin{abstract}
Motivated by the open problem of classifying Hopf $2$-cocycles for affine algebraic reductive groups $G$ over $\mathbb{C}$, in particular by \cite[Question 7.2]{EG1} which concerns minimal Hopf $2$-cocycles, we classify (minimal) Hopf $2$-cocycles for affine algebraic reductive groups $G$ whose connected component of the identity is a torus $T$. To achieve this we utilize \cite[Proposition 5.4]{ENO} in our situation to classify finite indecomposable semisimple module categories over the infinite semisimple equivariantization tensor category $\Rep(T)^K\simeq \Rep(G)$, where $K:=G/T$. We then use it to classify the module categories over $\Rep(G)$ of rank $1$, and show that the corresponding fiber functors on $\Rep(G)$ are classical (that is, preserve dimensions), which implies that they are in bijection with Hopf $2$-cocycles for $G$.
\end{abstract}

\maketitle

\tableofcontents

\section{Introduction}\label{sec:Introduction}
Let $G$ be an affine algebraic group over $\mathbb{C}$, and let $\Rep(G)$ be the tensor category of finite dimensional rational representations of $G$. Recall that {\bf Hopf $2$-cocycles} $J$ for $G$ (see \S\ref{sec:Hopf 2-cocycles}) are the same thing as tensor structures $J$ on the forgetful functor $F:\Rep(G)\to \Vect$. We will call such tensor functors $(F,J):\Rep(G)\to \Vect$ {\bf classical fiber functors} on $\Rep(G)$.

The classification of Hopf $2$-cocycles for {\bf finite} groups $G$ was achieved by Movshev in 
\cite{Mo}, and the classification for affine algebraic {\bf nilpotent} groups $G$ was achieved in \cite{EG1,EG2,G}. However, a general classification of Hopf $2$-cocycles for affine algebraic reductive groups $G$ is still unknown.

From now on, we let $G$ be an affine algebraic {\bf reductive} group over $\mathbb{C}$, so that its function algebra $\mathscr{O}(G)$ is a commutative {\bf cosemisimple} Hopf algebra (see e.g., \cite{A,M,S}). Since it is sufficient to classify {\bf minimal} Hopf $2$-cocycles for $G$ (see \S\ref{sec:cothopalg}), a positive answer to the following question will play an essential role in the classification of classical fiber functors on $\Rep(G)$.

\begin{question}{\rm \cite[Question 7.2]{EG1}}\label{redtor}
Suppose that an affine algebraic reductive  group $G$ admits a {\bf minimal} Hopf $2$-cocycle. Is it true that the connected component of the identity in $G$ is a torus? \qed
\end{question}

This paper is motivated by Question \ref{redtor}, and its main goal is to classify fiber functors on $\Rep(G)$ for an affine algebraic reductive group $G$ whose connected component of the identity is a {\bf torus} $T$. It turns out (see Theorem \ref{modcatr2}) that in this situation, all fiber functors $(F,J):\Rep(G)\to\Vect$ are classical; that is, $F$ is the forgetful functor. Thus, the classification of fiber functors on $\Rep(G)$ in our case is equivalent to the classification of Hopf $2$-cocycles for $G$.

To achieve our goal we first show in Theorem \ref{red2}
that there is an equivalence of tensor categories $\Rep(T)^K\simeq \Rep(G)$, where $K:=G/T$ (it is a finite group), and apply \cite[Proposition 5.4]{ENO} in our infinite situation to classify in Theorem \ref{partaunew} finite indecomposable semisimple module categories over $\Rep(T)^K$ in terms of $L$-equivariantizations of indecomposable semisimple module categories over ${\rm Rep}(T)$ for subgroups $L\subseteq K$ (see \S\ref{sec:Group actions on tensor categories}-\S\ref{sec:Equivariantization of module categories} for all the necessary definitions and constructions). We then use it to classify in Theorem \ref{modcatr2} the module categories over $\Rep(T)^K$ of rank $1$, hence fiber functors on $\Rep(G)$, and show that they are all classical. Finally, we use it to classify (minimal) Hopf $2$-cocycles for $G$ in Theorem \ref{fibredred1}. We conclude with the classification of (minimal) Hopf $2$-cocycles for direct product groups $G=T\times K$ in Corollary \ref{directproductcase}, and in the case where $K$ is also commutative, we compare our classification with the well known classification given by the Abelian group $H^2(\widehat{T\times K},\mathbb{C}^{\times})$ (see \S\ref{sec:appendix}).\\

\noindent
{\bf Acknowledgements.} This work was supported by Simons Foundation Award 963288.

\section{Preliminaries}\label{sec:Preliminaries}

\subsection{Hopf $2$-cocycles}\label{sec:Hopf 2-cocycles} Let $H$ be a Hopf algebra over $\mathbb{C}$. An element  
$J\in (H^{\ot 2})^*$ is called a {\bf Hopf $2$-cocycle} for $H$ if it is invertible under the convolution product, and satisfies
\begin{align*}
J (a_1b_1,c)J (a_2,b_2)= J
(a,b_1c_1)J (b_2,c_2),\quad J (a,1) =\varepsilon(a)=J (1,a)
\end{align*}
for every $a,b,c\in H$.

Given a Hopf $2$-cocycle $J$ for $H$, one can construct a new Hopf algebra $H^{J}$ as follows. As a coalgebra $H^{J}=H$, the new multiplication $m^{J}$ is given by
\begin{equation*}
m^{J}(a\ot b):=J^{-1} (a_1,b_1)a_2b_2J
(a_3,b_3),\quad a,b\in H,
\end{equation*}
and the new antipode is given by
\begin{equation*}
S^{J}(a)=J ^{-1}(a_1,S(a_2))S(a_3)J
(S(a_4),a_5),\quad a\in H.
\end{equation*} 
Equivalently, a Hopf $2$-cocycle $J$ for $H$ defines a tensor structure on the {\bf forgetful functor} $F:\Corep(H)\xrightarrow{}\Vect$, where $\Corep(H)$ is the category of finite dimensional right $H$-comodules. Namely, for every $X,Y\in \Corep(H)$, we have
\begin{equation*}
J_{X,Y}:X\ot Y\xrightarrow{\cong} X\ot Y,\,\,\,x\ot y\mapsto J^{-1}\cdot (x\ot y):=J^{-1}((x\ot y)^{(-1)})(x\ot y)^{(0)}.
\end{equation*}

Let $J$ be a Hopf $2$-cocycle for $H$. Recall that for any $x\in (H^*)^{\times}$ there is a Hopf $2$-cocycle $J^x$ for $H$, defined by 
$$ J^{x} (a,b):=x(a_1b_1) J(a_2, b_2) x^{-1}(a_3)x^{-1}(b_3),\quad a,b \in H,$$
and that two Hopf $2$-cocycles $J_1$ and $J_2$ for $H$ are {\bf gauge equivalent} if $J_2=J_1^x$ for some $x\in (H^*)^{\times}$. Recall also that gauge equivalence classes of Hopf $2$-cocycles for $H$ are in bijection with equivalence classes of tensor structures on the forgetful functor $F:\Corep(H)\to\Vect$. 

\subsection{Cotriangular Hopf algebras}\label{sec:cothopalg} 
Let $H$ be a Hopf algebra. Recall that $(H,R)$ is {\bf cotriangular} if $R$ is an invertible element of $(H^{\ot 2})^*$ that satisfies 
\begin{equation*}
R^{-1}=R_{21},\quad R(a,bc)=R(a_1,b)R(a_2,c),
\end{equation*}
\begin{equation*}
R(ab,c)=R(b,c_1)R(a,c_2),\quad R(a_1,b_1)b_2a_2=a_1b_1R(a_2,b_2)
\end{equation*}
for every $a,b,c\in H$.
Recall that $(H,R)$ is called {\bf minimal} if $R$ is nondegenerate. By \cite[Proposition 2.1]{G}, any cotriangular Hopf algebra $(H,R)$ has a unique minimal cotriangular Hopf algebra quotient.

Let $(H,R)$ be a cotriangular Hopf algebra, and let $J$ be a Hopf $2$-cocycle for $H$. Recall that $(H^J,R^J)$ is a cotriangular Hopf algebra, where $R^J:=J_{21}^{-1}RJ$, and that $J$ is called {\bf minimal} if $(H^J,R^J)$ is minimal; that is, if $R^J$ is nondegenerate.

\subsection{Group actions on semisimple tensor categories}\label{sec:Group actions on tensor categories}
Fix a {\bf finite} group $K$, and let ${\rm Cat}(K)$ be the
monoidal category whose objects are elements of $K$, morphisms are
identities, and tensor product is multiplication in
$K$. 

Fix a (not necessarily finite) {\bf semisimple} tensor category $\C$, and let $\text{Aut}^\ot(\C)$ denote the
monoidal category whose objects are tensor autoequivalences $a=(a,\alpha)$ of $\C$, morphisms are
natural isomorphisms, and tensor product is composition.
 
Recall (see, e.g., \cite[Definition 2.7.1]{EGNO}) that an {\bf action} of $K$ on $\C$ is a monoidal functor $(A,\gamma):{\rm Cat}(K) \to \text{Aut}^\ot(\C)$; that is,  
\begin{equation}\label{actiofunctor}
A:{\rm Cat}(K)\to \text{Aut}^\ot(\C),\,\,k \mapsto
(A_k,\alpha_k),
\end{equation}
is a functor equipped with a monoidal structure 
\begin{equation}\label{gammamonstr}
\gamma=\{\gamma_{k,\ell}: (A_k,\alpha_k)\circ (A_{\ell},\alpha_{\ell})\xrightarrow{\cong} (A_{k\ell},\alpha_{k\ell})\mid k,\ell\in K\}.
\end{equation}

\begin{example}\label{eqvec}
The group $K$ acts trivially on $\C:=\Vect $; that is, $A_k=\id_{\C}$, $\alpha_k=\id$, and $\gamma_{k,\ell}=\id_{{\rm Aut}^\ot(\C)}$ for every $k,\ell\in K$. \qed
\end{example}

\subsection{Equivariantization of semisimple tensor categories}\label{sec:Equivariantization of tensor categories}
Retain the notation from \S\ref{sec:Group actions on tensor categories}. Recall (see, e.g., \cite[Definition 2.7.2]{EGNO}) that a {\bf $K$-equivariant object} of $\C$ is a
pair $(X,u)$, where $X\in \C$ and 
$$u=\{u_k: A_k(X)\xrightarrow{\cong} X\mid k\in K\}$$ 
is a collection of $\C$-isomorphisms, such that the diagram
\begin{equation}
\label{equivariantX}
\xymatrix{A_k(A_{\ell}(X))\ar[rr]^{A_k(u_{\ell})} \ar[d]_{\gamma_{k,\ell}(X)
}&&A_k(X)\ar[d]^{u_k}\\ A_{k\ell}(X)\ar[rr]^{u_{k\ell}}&&X}
\end{equation}
commutes for all $k,\ell\in K$.
A {\bf morphism of equivariant objects} is a morphism in $\C$ which commutes with the equivariant structures. Equivariant objects together with their morphisms form a tensor category $\C^K$, called the {\bf $K$-equivariantization} of $\C$.

\begin{example}\label{eqvec1}
Keep the setup of Example \ref{eqvec}. In this case we see that a $K$-equivariant object $(X,u)$ in $\Vect$ is nothing but a (left) $K$-representation 
$$u:K\to {\rm GL}(X),\,\,\,k\mapsto u_k,$$ 
so $\Vect ^K\simeq \Rep(K)$ as tensor categories. \qed 
\end{example}

\subsection{Twisted module categories}\label{sec:Twisted module categories} 
Fix a left module category $(\M,m)$ over $\C$, with action $\overline{\ot}$. For every $(a,\alpha)\in \text{Aut}^\ot(\C)$, let $(\M,m)^{(a,\alpha)}$ denote the left $\C$-module category obtained from $(\M,m)$ by twisting the action of $\C$ by means of $(a,\alpha)$; that is,  
\begin{equation}\label{twistedM}
X\overline{\ot}^{(a,\alpha)} M:= a(X)\overline{\ot} M,\quad X\in\C,\,M\in \M,
\end{equation}
and for every $M\in \M$ and $X,Y\in \C$, the module associativity 
\begin{equation}\label{twistedmodstrM}
m^{(a,\alpha)}(X,Y,M):(X\ot Y)\overline{\ot}^{(a,\alpha)} M\xrightarrow{\cong}X\overline{\ot}^{(a,\alpha)}(Y\overline{\ot}^{(a,\alpha)} M)
\end{equation}
is given by the composition $m(a(X),a(Y),M)\circ (\alpha(X,Y)\ot \id_{M})$:
$$a(X\ot Y)\overline{\ot} M\xrightarrow{\alpha(X,Y)\ot \id_{M}}(a(X)\ot a(Y))\overline{\ot} M \xrightarrow{m(a(X),a(Y),M)}a(X)\overline{\ot}(a(Y)\overline{\ot}M).$$

Now given a $K$-action $(A,\gamma)$ on $\C$ as in \S\ref{sec:Group actions on tensor categories}, set  
\begin{equation}\label{M^k}
(\M ,m)^k:=(\M,m)^{(A_k,\alpha_k)},\quad k\in K.
\end{equation}
Note that $\gamma$ (\ref{gammamonstr}) induces a set of $\C$-module equivalences 
$$
\{(\id_{\M},\Gamma_{k,\ell}):((\M,m)^{\ell})^{k}\xrightarrow{\simeq}(\M,m)^{k\ell}\mid k,\ell\in K\},
$$
where for every $X\in\C$ and $M\in \M$, we define 
$$\Gamma_{k,\ell}(X,M):=\gamma_{k,\ell}(X)\ot {\rm id}_M,$$
so that
\begin{equation}\label{Gamma}
\Gamma_{k,\ell}(X,M):X(\overline{\ot}^{\ell})^k M=A_{k}(A_{\ell}(X))\overline{\ot}M\xrightarrow{\cong}A_{k\ell}(X)\overline{\ot}M=X\overline{\ot}^{k\ell} M.
\end{equation}  
In particular, $K$ acts on the set of indecomposable exact $\C$-module categories.

\begin{example}\label{eqvec2}
Keep the setup of Example \ref{eqvec1}. Let $(\Vect,\id)$ be the unique indecomposable module category over $\Vect$. Then for every $k,\ell\in K$, we have 
$(\Vect,\id)^k=(\Vect,\id)$, and 
$\Gamma_{k,\ell}=\id$. \qed 
\end{example}

\subsection{Equivariant module categories}\label{sec:Equivariant module categories}
Retain the notation from \S\ref{sec:Twisted module categories}. Fix a {\bf subgroup} $L\subseteq K$. Recall \cite{ENO} that an {\bf $L$-equivariant structure} on $(\M,m)$ is a triple $(U,\tau,\mu)$, such that the following hold:

(1) $(U,\tau)=\{(U_{\ell},\tau_{\ell}):(\M,m)\xrightarrow{\simeq} (\M,m)^{\ell}\mid \ell\in L\}$ 
is a set of $\C$-module equivalences. Namely, for every $\ell\in L$,
$$U_{\ell}:\M\xrightarrow{\cong} \M$$
is an equivalence of Abelian categories, and 
$$\tau_{\ell}=\{\tau_{\ell}(X,M):U_{\ell}(X\overline{\ot}M)\xrightarrow{\cong} X\overline{\ot}^{\ell}U_{\ell}(M)\mid X\in\C,M\in\M\}$$
is a natural isomorphism, such that for every $X,Y\in\C$ and $M\in\M$, we have
\begin{equation}\label{tauass}
\begin{split}
& (\id_X\ot \tau_{\ell}(Y,M))\tau_{\ell}(X,Y\overline{\ot}M) U_{\ell}(m(X,Y,M))= m^{\ell}(X,Y,U_{\ell}(M))\tau_{\ell}(X\ot Y,M)\\
& \qquad \qquad\qquad\quad:U_{\ell}((X\ot Y)\overline{\ot} M)\xrightarrow{\cong} X\overline{\ot}^{\ell}(Y\overline{\ot}^{\ell} U_{\ell}(M)),
\end{split}
\end{equation}
and
\begin{equation}\label{tauunit}
l_{U_{\ell}(M)}\circ\tau_{\ell}(\mathbf{1},M)=U_{\ell}(l_M):U_{\ell}(M)\xrightarrow{\cong} U_{\ell}(M).
\end{equation}

Recall that for every $k\in L$, 
$$(U,\tau)^k=\{(U_{\ell}^k,\tau_{\ell}^k):=(U_{\ell},\tau_{\ell})^k:
(\M,m)^k\xrightarrow{\simeq}((\M,m)^{\ell})^k \mid \ell\in L\},$$
where for every $k,\ell\in L$,
$$(U_{\ell}^k,\tau_{\ell}^k):(\M,m)^k\xrightarrow{\simeq}
((\M,m)^{\ell})^k$$
is the $\C$-module equivalence with
$$U_{\ell}^k=U_{\ell}:\M\xrightarrow{\simeq} \M$$
as Abelian functors, and 
$$\tau_{\ell}^k(X,M)=\tau_{\ell}(X^k,M),\quad X\in\C,\,M\in\M.$$

(2) 
\begin{equation}\label{mufirst}
\mu=\{\mu_{k,\ell}:(U_{k},\tau_{k})\circ(U_{\ell},\tau_{\ell})^k \xrightarrow{\cong} (U_{k\ell},\tau_{k\ell})\circ(\id_{\M},\Gamma_{k,\ell})\mid k,\ell\in L\}
\end{equation} 
is a set of natural isomorphisms of $\C$-module functors $(\M,m)\to((\M,m)^{\ell})^k$, satisfying the following compatibility conditions:
$$(\mu_{i,k}(U_{\ell})^{ik})\circ (\mu_{ik,{\ell}}\Gamma_{i,k})=(U_i(\mu_{k,{\ell}})^i)\circ (\mu_{i,k{\ell}}\Gamma_{k,{\ell}}).$$
Namely, 
for every $k,\ell\in L$, we have a collection of $\M$-isomorphisms 
$$\{\mu_{k,\ell}(M): U_k(U_{\ell}(M))\xrightarrow{\cong}U_{k\ell}(M)\mid M\in \M\},$$
such that for every $X\in \C$, $M\in \M$ and $i,k,\ell\in L$, the following hold: 
\begin{equation}\label{munat0}
\tau_{k\ell}(X,M)\mu_{k,\ell}(X\overline{\ot} M)= (\gamma_{k,\ell}(X)\ot \mu_{k,\ell}(M))\tau_{k}(X^\ell,U_{\ell}\left(M\right))U_{k}(\tau_{\ell}(X,M));
\end{equation}
that is, the diagram
\begin{equation*}
\xymatrixcolsep{4.5pc}\xymatrix{
U_k(U_{\ell}(X\overline{\ot}M)) \ar[d]_{\mu_{k,\ell}(X\overline{\ot} M)} \ar[r]^{U_{k}(\tau_{\ell}(X,M))} 
& U_{k}(X\overline{\ot}^\ell U_{\ell}(M)) \ar[r]^{\tau_{k}(X^\ell,U_{\ell}\left(M\right))}
& X^{\ell}\overline{\ot}^{k} U_{k}(U_{\ell}(M)) \ar[d]^{\gamma_{k,\ell}(X)\ot \mu_{k,\ell}(M)}\\
U_{k\ell}(X\overline{\ot}M) \ar[r]^{\tau_{k\ell}(X,M)} 
& X\overline{\ot}^{k\ell} U_{k\ell}(M) \ar[r]^{\id} 
& X\overline{\ot}^{k\ell}U_{k\ell}(M)
}
\end{equation*}
commutes, and 
\begin{equation}\label{comp00}
\mu_{i,k\ell}(M)U_i(\mu_{k,\ell}(M))=
\mu_{ik,\ell}(M)\mu_{i,k}(U_{\ell}(M)): U_{i}(U_k(U_{\ell}(M)))\xrightarrow{\cong}U_{ik\ell}(M).
\end{equation}

\begin{example}\label{eqvec3}
Keep the setup of Example \ref{eqvec2}. An $L$-equivariant structure on $(\Vect,\id)$ is given by 
$$U=\id,\,\,\tau=\id,\,\,\mu_{k,\ell}:\id_{\Vect}\xrightarrow{\cong} \id_{\Vect},\quad k,\ell\in L.$$ By (\ref{comp00}), $\mu_{k,\ell}$ is determined by the nonzero scalar $\mu_{k,\ell}(\mathbb{C})\cdot {\rm id}:\mathbb{C}\xrightarrow{\cong} \mathbb{C}$, and the corresponding function $\mu:L\times L\to\mathbb{C}^{\times}$, $(k,\ell)\mapsto \mu_{k,\ell}(\mathbb{C})$, is a $2$-cocycle. Thus, $L$-equivariant structures on $(\Vect,\id)$ over $\Vect$ correspond to elements of the group $Z^2(L,\mathbb{C}^{\times})$. \qed
\end{example}

\begin{example}\label{L=1}
Let $L=1$. Every $(\M,m)$ has a unique $L$-equivariant structure; the trivial one. \qed
\end{example}

\subsection{Equivariantization of module categories}\label{sec:Equivariantization of module categories}
Retain the notation from \S\ref{sec:Equivariant module categories}. Fix an $L$-equivariant structure $(U,\tau,\mu)$ on $(\M,m)$. An {\bf equivariant object} in $(\M,m)$ with respect to this structure is a pair $(M,v)$ consisting of an object $M\in \M$ along with a set of $\M$-isomorphisms 
$$v=\{v_{\ell}: U_{\ell}(M)\xrightarrow{\cong} M\mid {\ell}\in L\},$$ such that the diagram
\begin{equation}\label{diagramv}
\xymatrix{U_k(U_{\ell}(M))\ar[d]_{\mu_{k,\ell}(M)}\ar[rr]^{U_k(v_{\ell})}&& U_k(M)\ar[d]^{v_k} &&\\
U_{k\ell}(M)\ar[rr]^{v_{k\ell}}&& M&&}
\end{equation}
commutes for all $k,\ell\in L$. The {\bf $L$-equivariantization} category
$$((\M,m),U,\tau,\mu)^L$$ 
is the category whose objects are $L$-equivariant objects in $(\M,m)$, and morphisms $f:(M,v)\to (N,w)$ are morphisms $f:M\to N$ in $\M$, such that the diagram
\begin{equation}\label{diagramvmor}
\xymatrix{U_{\ell}(M)\ar[d]_{U_{\ell}(f)}\ar[rr]^{v_{\ell}}&& M\ar[d]^{f} &&\\
U_{\ell}(N)\ar[rr]^{w_{\ell}}&& N&&}
\end{equation}
commutes for all $\ell\in L$.

\begin{example}\label{eqvec4}
Keep the setup of Example \ref{eqvec3}. In this case we have 
$$((\Vect,\id),\id,\id,\mu)^L={\rm Rep}(L,\mu^{-1});$$ 
that is, the category of finite dimensional projective representations of $L$ with the $2$-cocycle $\mu^{-1}$. \qed
\end{example}

Recall \cite{ENO} that the category $((\M,m),(U,\tau),\mu)^L$ is a left module category over $\C^K$ as follows.
For every $(M,v)\in ((\M,m),(U,\tau),\mu)^L$ and $(X,u),(Y,w)\in \C^K$, the action $\ot^L$ of $\C^K$ on $((\M,m),(U,\tau),\mu)^L$ is given by $$(X,u)\overline{\ot}^{L}(M,v):=(X\overline{\ot} M,u\overline{\ot} v),$$ where,   
$$u\overline{\ot}v:=\{(u\overline{\ot}v)_{\ell}:=(u_{\ell}\overline{\ot} v_{\ell})\circ\tau_{\ell}(X,M):U_{\ell}(X\overline{\ot}M)\xrightarrow{\cong} X\overline{\ot}M\mid \ell\in L\},$$ and the module associativity constraint
$$m^L(X,Y,M):((X\ot Y)\overline{\ot} M,(u\ot w)\overline{\ot}v))\xrightarrow{\cong}(X\overline{\ot}(Y\overline{\ot}M),u\overline{\ot}(w\overline{\ot}v))
$$
is given by the module associativity constraint $m(X,Y,M)$ of $\M$.

\begin{example}\label{eqvec5}
Keep the setup of Example \ref{eqvec4}. Then indecomposable semisimple module categories over $\Rep(K)=\Vect^K$ are classified by conjugacy classes of pairs $(L,\mu)$, $(L,\mu)\mapsto ((\Vect,\id),\id,\id,\mu)^L ={\rm Rep}(L,\mu^{-1})$, where $L\subseteq K$ is a subgroup and $\mu\in Z^2(L,\mathbb{C}^{\times})$. 

In particular, the module category $\V:=((\Vect,\id),\id,\id,\mu)^L={\rm Rep}(L,\mu^{-1})$ has rank $1$ if and only if 
$\mu$ is {\bf nondegenerate}, so $L$ is a group of central type. In this case, let $(V,v)$ be the unique simple object of $\V={\rm Rep}(L,\mu^{-1})$.
Then for every $(X,u)\in \Rep(K)=\Vect^K$, we have
$$(X,u)\ot ^{\V}(V,v)=(X\ot V,u\ot v)\in {\rm Rep}(L,\mu^{-1}),$$
where  
$$u\ot v=\left\{\left(u\ot v\right)_{\ell}=u_{\ell}\ot v_{\ell}:X\ot V\xrightarrow{\cong}X\ot V\mid \ell\in L\right\},$$ 
and the module associativity constraint in $\V$ is the standard associativity constraint of vector spaces. Indeed, 
the diagram
\begin{equation}\label{diagramv1}
\xymatrix{X\ot V\ar[d]_{\mu_{k,\ell}}\ar[rr]^{u_{\ell}\ot v_{\ell}}&& X\ot V\ar[d]^{u_{k}\ot v_{k}} &&\\
X\ot V\ar[rr]^{u_{k\ell }\ot v_{k\ell}}&& X\ot V&&}
\end{equation}
commutes.

Finally, it is clear that $\V$ corresponds to a minimal Hopf $2$-cocycle for $K$ if and only if $L=K$; that is, if and only if $K$ is a group of central type with nondegenerate $2$-cocycle $\mu\in Z^2(K,\mathbb{C}^{\times})$. \qed
\end{example}
 
\begin{proposition}\label{rankddec}
%If $\C$ and $\M$ are semisimple, then  
For any finite semisimple $\M:=((\M,m),U,\tau,\mu)$ and subgroup $L\subseteq K$, we have that ${\rm rank}(\M^L)\ge {\rm rank}(\M)$. 
In particular, if $((\M,m),U,\tau,\mu)^L$ has rank $1$, then $\M$ is a $\C$-module category of rank $1$.
\end{proposition}

\begin{proof}
Since the forgetful functor $F:\M^L\to \M$, $(M,v)\mapsto M$, has an adjoint $I:\M\to \M^L$, it follows that every simple $M\in \M$ is a direct summand of $F(I(M))$ since $\Hom_{\M}(M,F(I(M)))=\Hom_{\M^L}(I(M),I(M))$. Thus, $F$ is surjective, which implies that ${\rm rank}(\M^L)\ge {\rm rank}(\M)$. The second claim follows easily from the first one.
\end{proof}

Consider now the $K$-graded crossed product tensor category $\C\rtimes K$ (see, e.g., \cite[Definition 4.15.5]{EGNO}), and recall that $\C$ admits a natural structure of an indecomposable semisimple left module category over $\C\rtimes K$, given by
$$(X\boxtimes k)\ot Y:=A_k(X\ot Y),\quad X,Y\in \C,\,k\in K.$$
Recall also (see, e.g., \cite[Definition 7.12.2]{EGNO}) the dual tensor category 
\begin{equation}\label{dtc}
(\C\rtimes K)^*_{\C}:={\rm Fun}_{\C\rtimes K}(\C,\C).
\end{equation}
Let ${\rm Mod}(\C\rtimes K)$ and ${\rm Mod}((\C\rtimes K)^*_{\C})$ be the $2$-categories of finite semisimple module categories over $\C\rtimes K$ and $(\C\rtimes K)^*_{\C}$, respectively. Since $\Fun_{\C\rtimes K}(\C,\mathscr{M})$ and $\Fun_{(\C\rtimes K)^*_{\C}}(\C,\mathscr{N})$ are finite for every $\mathscr{M}\in {\rm Mod}(\C\rtimes K)$ and $\mathscr{N}\in {\rm Mod}((\C\rtimes K)^*_{\C})$, we have well defined $2$-functors
$${\rm Mod}(\C\rtimes K)\to {\rm Mod}((\C\rtimes K)^*_{\C}),\,\,
\mathscr{M}\mapsto\Fun_{\C\rtimes K}(\C,\mathscr{M}),$$ 
and
$${\rm Mod}((\C\rtimes K)^*_{\C})\to {\rm Mod}(\C\rtimes K),\,\,
\mathscr{N}\mapsto\Fun_{(\C\rtimes K)^*_{\C}}(\C,\mathscr{N}).$$

\begin{lemma}\label{implem}
%Let $\C$ be a {\bf semisimple} tensor category, with an action of a finite group $K$. 
The following hold:
\begin{enumerate}
\item
There is an equivalence of tensor categories $(\C\rtimes K)^*_{\C}\simeq \C^K$.
\item
The above $2$-functors are inverse to each other.
Thus, the finite indecomposable semisimple module categories over $\C^K$ are in bijection with those over $\C\rtimes K$.
\end{enumerate}
\end{lemma}

\begin{proof}
(1) This is similar to the fusion case (see, e.g., \cite[Example 7.12.25]{EGNO}).

(2) This follows from (1) in exactly the same way it does in the fusion case.
\end{proof}

\begin{theorem}\label{eno}
Every finite indecomposable semisimple $\C^K$-module category is equivalent to one of the form $((\M,m),U,\tau,\mu)^L$ for some subgroup $L\subseteq K$ and an 
indecomposable semisimple $L$-equivariant $\C$-module category $((\M,m),U,\tau,\mu)$.
\end{theorem}

\begin{proof}
The proof is similar to \cite[Proposition 5.4]{ENO} using Lemma \ref{implem}. 
\end{proof}

\subsection{Classical fiber functors on ${\rm Corep}(H)^K$}\label{sec:Fiber functors}
Let $\C:={\rm Corep}(H)$ be the semisimple tensor category of finite dimensional comodules over a {\bf cosemisimple} Hopf algebra $H$. Let $J$ be a Hopf $2$-cocycle for $H$. Then $(\Vect,J)$ 
is a $\C$-module category of rank $1$, or equivalently, $(F,J):{\rm Corep}(H)\to \Vect$ is a classical fiber functor, so that $F(X)=\overline{X}$. 

Recall that ${\rm Aut}(F)\cong (H^*)^{\times}$ as groups.

\begin{theorem}\label{partaunewgen}
Fix a subgroup $L\subseteq K$. 
Then $L$-equivariant structures on $(\Vect,J)$ over ${\rm Corep}(H)$ are parametrized by pairs $(\tau,\mu)$, where 
\begin{equation}\label{taufirsttgen}
\tau: L\to (H^*)^{\times},\,\,\,\ell\mapsto \tau_{\ell},
\end{equation}
and
\begin{equation}\label{mufirstt0gen}
\mu\in Z^2(L,\mathbb{C}^{\times}),\,\,\,(k,\ell)\mapsto \mu_{k,\ell},
\end{equation}
such that for every $k,\ell\in L$ the following hold:
\begin{equation}\label{taulPsigen}
(\tau_{\ell}\ot\tau_{\ell})J= J^{\ell}\Delta(\tau_{\ell}),\,\,\,\,
\tau_{\ell}(\mathbf{1})=1,
\end{equation}
and
\begin{equation}\label{Ultaulgen}
\tau_k\tau_{\ell}^k=\tau_{k\ell}\gamma_{k,\ell}.
\end{equation}

In particular, if $(\Vect,J)$ admits an $L$-equivariant structure over ${\rm Corep}(H)$, then $J$ and $J^{\ell}$ are gauge equivalent Hopf $2$-cocycles for $H$ for every $\ell\in L$. 
\end{theorem}

\begin{proof}
Suppose that $(U,\tau,\mu)$ is an $L$-equivariant structure on $(\Vect,J)$ over ${\rm Corep}(H)$. Recall \S\ref{sec:Equivariant module categories} that this means the following:

(1) $(U,\tau)=\{(U_{\ell},\tau_{\ell}):(\Vect,J)\xrightarrow{\simeq} (\Vect,J^{\ell})\mid \ell\in L\}$ 
is a set of ${\rm Corep}(H)$-module equivalences. Then for every $\ell\in L$, we have that $U_{\ell}=\id$, and   
$$\tau_{\ell}=\{\tau_{\ell}(X):=\tau_{\ell}(X,\mathbb{C}):\overline{X}\xrightarrow{\cong}\overline{X^{\ell}}=\overline{X}\mid X\in {\rm Corep}(H)\}$$
is a natural isomorphism. Thus, for every $\ell\in L$, $\tau_{\ell}\in {\rm Aut}(F)=(H^*)^{\times}$, which gives  (\ref{taufirsttgen}). 

Moreover, it follows from (\ref{tauass})-(\ref{tauunit}) that for every $X,Z\in {\rm Corep}(H)$, we have
\begin{gather*}
(\id_{X}\ot \tau_{\ell}(Z))\tau_{\ell}(X,\overline{Z})m_{J}(X,Z,\mathbb{C})
= m_{J}^{\ell}(X,Z,\mathbb{C})\tau_{\ell}(X\ot Z)\\
:(\overline{X}\ot \overline{Z})\overline{\ot} \mathbb{C}\xrightarrow{\cong} X\overline{\ot}^{\ell}(Z\overline{\ot}^{\ell} \mathbb{C}),
\end{gather*}
and
\begin{equation*}
l_{\mathbb{C}}\circ\tau_{\ell}(\mathbf{1})=l:\mathbb{C}\xrightarrow{\cong}\mathbb{C}.
\end{equation*}
Thus, for every $\ell\in L$, we have
$$\tau_{\ell}(Z)\tau_{\ell}(X)J(X,Z)= J^{\ell}(X,Z)\tau_{\ell}(X\ot Z),\quad \tau_{\ell}(\mathbf{1})=1,$$
which yield (\ref{taulPsigen}).

(2) $\mu=\{\mu_{k,\ell}: (\id,\tau_{k})\circ(\id,\tau_{\ell})^k\xrightarrow{\cong}(\id,\tau_{k\ell})\circ(\id_{\M},\Gamma_{k,\ell})\mid k,\ell\in L\}$ is a set of natural isomorphisms of module functors $(F,J)\to (F,J^{k\ell})$. Thus, 
for every $k,\ell\in L$, we have a nonzero scalar 
$$\mu_{k,\ell}:=\mu_{k,\ell}(\mathbb{C}): \mathbb{C}\xrightarrow{\cong}\mathbb{C},$$
so (\ref{mufirstt0gen}) holds. Then by (\ref{munat0}), we have 
$$
\tau_k(X)\tau_{\ell}(X^k)\mu_{k,\ell}=\mu_{k,\ell}\tau_{k\ell}(X)\Gamma_{k,\ell}(X)$$
for every $X\in {\rm Corep}(H)$ and $k,\ell\in L$,  
and by (\ref{comp00}), we have 
$$
\mu_{i,k}\mu_{ik,\ell}=\mu_{k,\ell}\mu_{i,k\ell}
$$
for every $i,k,\ell\in L$. 
Thus, (\ref{mufirstt0gen}) and (\ref{Ultaulgen}) hold.

Conversely, it is clear that any pair $(\tau,\mu)$ satisfying (\ref{taufirsttgen})-(\ref{Ultaulgen}) defines an $L$-equivariant structure on $(\Vect,J)$ over ${\rm Corep}(H)$.
\end{proof}

\section{Hopf $2$-cocycles for certain reductive groups}\label{sec:The reductive case} 

%In this section we apply Theorem \ref{partaunewgen} to generalize Example \ref{eqvec5} to the following setting.
In this section we assume that $G$ is an affine algebraic reductive group whose connected component of the identity element is a torus $T$, and set $K:=G/T$. Then $\Rep(G)={\rm Corep}(\mathscr{O}(G))$ is semisimple tensor category.

\begin{example}
A typical example of such a group $G$ is the following. Let $\tilde{G}$ be an affine algebraic reductive group with maximal torus $T$, and take $G\subseteq \tilde{G}$ to be the normalizer of $T$.
%
%It is known that if $K:=G/T$ is a finite $p$-group, then the dimension of any irreducible representation of $G$ is a power of $p$. 
\qed
\end{example}

First we show that $\Rep(G)\simeq\Rep(T)^K$ as tensor categories. Second we classify 
finite semisimple module categories over $\Rep(T)^K$, and then deduce from it a classification of rank-one module categories
(fiber functors) over $\Rep(G)$. Finally we prove that every fiber functor on $\Rep(G)$ is classical, and obtain a classification of (minimal) Hopf $2$-cocycles for $G$.

\subsection{The tensor equivalence $\Rep(G)\simeq\Rep(T)^K$}\label{sec:The tensor equivalence}
Consider the exact sequence
\begin{equation}\label{sesalgr}
1\to T\to G\xrightarrow{\pi}K\to 1.
\end{equation}
Fix a section $\sigma:K\xrightarrow{1:1} G$ of $\pi$ with $\sigma(1)=1$.

\begin{lemma}\label{cohclassgamma}
The following hold:
\begin{enumerate}
\item
For every $k\in K$ and $t\in T$, the element $t^k:=\sigma(k)t\sigma(k)^{-1}\in T$ does not depend on the choice of $\sigma$.
\item
For every $k,\ell\in K$, $\eta(k,\ell):=\sigma(k)\sigma(\ell)\sigma(k\ell)^{-1}\in T$. In particular, $K$ acts on $T$ from the left via $k\cdot t:=t^k$, $k\in K$ and $t\in T$.
\item
The function
\begin{equation*}
\eta:K\times K\to T,\,\,\,(k,\ell)\mapsto \eta(k,\ell), 
\end{equation*}
is a $2$-cocycle in $Z^2(K,T)$, whose class $[\eta]\in H^2(K,T)$ does not depend on the choice of $\sigma$.
\item
Let $N:=\{t\in T\mid t^{|K|}=1\}\subset T$. Then $[\eta]\in H^2(K,N)\subset H^2(K,T)$.
\item
As varieties, $G=T\times K$, and the product in $G$ is given by
\begin{equation*}
(t,{k})\cdot(s,{\ell})=\left(ts^k\eta(k,\ell),{k}{\ell}\right),\quad s,t\in T,\,k,\ell\in K.
\end{equation*}
\end{enumerate}
\end{lemma}

\begin{proof}
(1) This follows since $T$ is commutative.

(2) Since $\pi(\eta(k,\ell))=\pi(\sigma(k)\sigma(\ell)\sigma(k\ell)^{-1})=k\ell(k\ell)^{-1}=1$, we have $\eta(k,\ell)\in T$ for every $k,\ell\in K$. Also, for every $k,\ell\in K$ and $t\in T$, we have
\begin{equation*}
\begin{split}
& (k\ell)\cdot t=\sigma(k\ell)t\sigma(k\ell)^{-1}=\left(\eta(k,\ell)^{-1}\sigma(k)\sigma(\ell)\right)t\left(\eta(k,\ell)^{-1}\sigma(k)\sigma(\ell)\right)^{-1}\\
&=\left(\eta(k,\ell)^{-1}\sigma(k)\sigma(\ell)\right)t\left(\sigma(\ell)^{-1}\sigma(k)^{-1} \eta(k,\ell)\right)\\
&=\left(\eta(k,\ell)^{-1}\sigma(k)\right)(\ell \cdot t)\left(\sigma(k)^{-1} \eta(k,\ell)\right)\\
&=\eta(k,\ell)^{-1}\left(k\cdot (\ell \cdot t)\right)\eta(k,\ell)=k\cdot (\ell \cdot t),
\end{split}
\end{equation*}
where the last equality follows since both $\eta(k,\ell)$ and $k\cdot (\ell \cdot t)$ lie in $T$.

(3) We have to verify that for every $k,\ell,m\in K$, we have
$$(k\cdot \eta(\ell,m))\eta(k,\ell m)=\eta(k\ell,m)\eta(k,\ell).$$
Indeed, on the one hand, we have
\begin{equation*}
\begin{split}
& \eta(k,\ell m)\eta(k\ell,m)^{-1}\\
& = \sigma(k)\sigma(\ell m)\sigma(k\ell m)^{-1}(\sigma(k\ell)\sigma(m)\sigma(k\ell m)^{-1})^{-1}\\
& = \sigma(k)\sigma(\ell m)\sigma(k\ell m)^{-1}\sigma(k\ell m)\sigma(m)^{-1}\sigma(k\ell)^{-1}\\
& = \sigma(k)\sigma(\ell m)\sigma(m)^{-1}\sigma(k\ell)^{-1},
\end{split}
\end{equation*}
and on the other hand, we have
\begin{equation*}
\begin{split}
& \eta(k,\ell)(k\cdot \eta(\ell,m))^{-1}=\eta(k,\ell)(\eta(\ell,m)^k)^{-1}\\
& = \left(\sigma(k)\sigma(\ell)\sigma(k\ell)^{-1}\right)\left(\sigma(k)\eta(\ell,m)\sigma(k)^{-1}\right)^{-1}\\
& = \left(\sigma(k)\sigma(\ell)\sigma(k\ell)^{-1}\right)\left(\sigma(k)\eta(\ell,m)^{-1}\sigma(k)^{-1}\right)\\
& = \left(\sigma(k)\sigma(\ell)\sigma(k\ell)^{-1}\right)\left(\sigma(k)(\sigma(\ell)\sigma(m)\sigma(\ell m)^{-1})^{-1}\sigma(k)^{-1}\right)\\
& =  \left(\sigma(k)\sigma(\ell)\sigma(k\ell)^{-1}\right)\left(\sigma(k)\sigma(\ell m)\sigma(m)^{-1}\sigma(\ell)^{-1}\sigma(k)^{-1}\right)\\
& =  \left(\sigma(k)\sigma(\ell m)\sigma(m)^{-1}\sigma(\ell)^{-1}\sigma(k)^{-1}\right) \left(\sigma(k)\sigma(\ell)\sigma(k\ell)^{-1}\right)\\
& = \sigma(k)\sigma(\ell m)\sigma(m)^{-1}\sigma(k\ell)^{-1},
\end{split}
\end{equation*}
where the last equality follows since $T$ is commutative.

(4) This follows from $|K|H^2(K,T)=0$.

(5) This follows in a standard way.
\end{proof}

Let $\C:={\rm Rep}(T)=\Vect(\widehat{T})$, where $\widehat{T}$ is the group of characters of $T$. Then $\C$ is a semisimple  pointed tensor category with simple objects $\{\delta_x\mid x\in \widehat{T}\}$.

For every $X\in \Rep(T)$ and $ k\in K$, let $X^{ k}\in \Rep(T)$ be the $\Ad(k)$-twisted representation of $T$ on $\overline{X}$ (= the underlying vector space of $X$); that is, the action of $t\in T$ on $\overline{X}$ is given by the action of $t^k=ktk^{-1}\in T$. Using the $2$-cocycle $\eta$ from Lemma \ref{cohclassgamma} we can define an action $(A,\gamma)$ of $K$ on $\Rep(T)$ as follows: 
\begin{equation}\label{actionnew}
A_{{k}}(X)=\, X^{{k}}\,\,\,\text{and}\,\,\,\gamma_{{k},{\ell}}(X):=\eta(k,\ell)_{|X}:(X^{{\ell}})^{{k}}\xrightarrow{\cong} X^{{k\ell}},
\end{equation}
for every $X\in \Rep(T)$ and $k,\ell\in K$ (the tensor structure $\alpha_k$ on $A_k$ being the identity for every $k\in K$). Indeed, the $2$-cocycle condition on $\eta$ implies that 
$$\gamma:=\{\gamma_{k,\ell}:A_k\circ A_{\ell}\xrightarrow{\cong}A_{k\ell}\mid k,\ell\in K\}$$
defines a monoidal structure on $A$. Note that each $\gamma_{k,\ell}$ is determined by the set of nonzero scalars  
$$
\{\gamma_{{k},{\ell}}(\delta_{x})=x(\eta(k,\ell))\mid x\in\widehat{T}\}\subset\mathbb{C}^{\times}.
$$

Now let ${\rm X}:=(X,u)\in \Rep(T)^K$, so that $X\in \Rep(T)$ and 
$$u=\{u_{{k}}:A_{{k}}(X)=X^{{k}}\xrightarrow{\cong} X\mid k\in K\}$$ 
is a set of $T$-isomorphisms satisfying (\ref{equivariantX}); that is, for every $k,\ell\in K$, we have 
$$u_{{k}{\ell}}\circ \gamma_{{k},{\ell}}(X)=u_{{k}}\circ A_k(u_{{\ell}})=u_ku_{\ell}$$ 
in $\GL(\overline{X})$. Thus, we have a representation $\rho_{(X,u)}$ of $G$ on $\overline{X}$, given by
$$\rho_{(X,u)}:G\to \GL(\overline{X}),\,\,(t,{k})\mapsto t_{|X}\circ u_{{k}},$$
so we have determined a tensor functor
\begin{equation}\label{fun}
\Rep(T)^K\to \Rep(G),\,\,(X,u)\mapsto \rho_{(X,u)}.
\end{equation}

Given a simple object $\delta_x\in\Rep(T)$, let $K_x\subseteq K$ be the stabilizer of $\delta_x$, let $\{g_1,\dots,g_r\}$ be a set of representatives of the double cosets of $K_x$ in $K$, let $\widehat{T}_x\subseteq \widehat{T}$ be the $K$-orbit (= $G$-orbit) of $\delta_x$, and let $G_x:=\pi^{-1}(K_x)\subseteq G$.

\begin{theorem}\label{red2}
The following hold:
\begin{enumerate}
\item
The tensor functor (\ref{fun}) is an equivalence $\Rep(T)^K\simeq \Rep(G)$.
\item
Isomorphism classes of simples of ${\rm Rep}(G)$ are in bijection with equivalence classes of pairs $(\widehat{T}_x,W_x)$, where $W_x\in\Rep(K_x,\eta)$ is simple. 
Explicitly, the simple $G$-module corresponding to $(\widehat{T}_x,W_x)$ is given by ${\rm Ind}_{G_x}^G(W_x\ot \delta_x)$, with forgetful image 
$\bigoplus_{i=1}^r W_x\ot_{\mathbb{C}}\delta_{x^{g_i}}\in\Rep(T)$.
\end{enumerate}
\end{theorem}

\begin{proof}
(1) To construct the inverse functor $\Rep(G)\to \Rep(T)^K$ take ${\rm X}\in \Rep(G)$, and consider its restriction $X:={\rm X}_{|T}$ to $T$.  
The action of $K$ on $T$ defines a set of $T$-isomorphisms $u:=\{u_k:A_k(X)=X^{k}\xrightarrow{\cong} X\mid k\in K\}$, where $u_k(x)=k^{-1}\cdot x$ (where we view $K$ as a subset of $G$). Then 
$(X,u)$ is a $K$-equivariant object of $\Rep(T)$, and it is clear that the tensor functor
$$\Rep(G)\to \Rep(T)^K,\,\,{\rm X}\mapsto ({\rm X}_{|T},u),$$
is inverse to (\ref{fun}).

(2) This follows from (1) in a standard way.
\end{proof}

\subsection{Finite semisimple module categories over $\Rep(G)$}\label{sec:Finite semisimple module categories}
Let $\widehat{S}\subseteq \widehat{T}$ be a subgroup of {\bf finite index}, and let $\Psi\in H^2(\widehat{T},C^1(\widehat{T}/\widehat{S},\mathbb{C}^{\times}))$. Let $\M(\widehat{S},\Psi)$ be the finite semisimple $\Rep(T)$-module category with simple objects 
$$\{\delta_{[y]}\mid [y]\in \widehat{T}/\widehat{S}\},$$ 
where $[y]:=y\widehat{S}$, $\Rep(T)$-action determined by
$$\delta_{x}\overline{\ot}\delta_{[y]}:=\delta_{[xy]},\quad x\in \widehat{T},\,\,[y]\in \widehat{T}/\widehat{S},$$
and module associativity constraint $m_{\Psi}$ determined by
$$m_{\Psi}(\delta_x,\delta_z,\delta_{[y]}):=\Psi(x,z)([y])\cdot \id:\delta_{[xzy]}=(\delta_{x}\ot \delta_{z})\overline{\ot}\delta_{[y]}\xrightarrow{\cong}\delta_{x}\overline{\ot}(\delta_{z}\overline{\ot}\delta_{[y]})=\delta_{[xzy]}$$
for every $x,z\in \widehat{T}$ and $[y]\in \widehat{T}/\widehat{S}$.

\begin{proposition}\cite{EGNO,G2}\label{fibtor}
Finite indecomposable semisimple module categories over $\Rep(T)$ are in bijection with 
pairs $(\widehat{S},\Psi)$, $\M(\widehat{S},\Psi)\mapsfrom (\widehat{S},\Psi)$, where $\widehat{S}\subseteq \widehat{T}$ is a subgroup of finite index, and $\Psi\in H^2(\widehat{T},C^1(\widehat{T}/\widehat{S},\mathbb{C}^{\times}))$. \qed
\end{proposition}

\begin{remark}\label{shapiro}
By Shapiro's lemma, we have a group isomorphism
$$H^2(\widehat{T},C^1(\widehat{T}/\widehat{S},\mathbb{C}^{\times}))\xrightarrow{\cong}H^2(\widehat{S},\mathbb{C}^{\times}),\,\,\,\Psi\mapsto \psi:=\Psi_{\mid \widehat{S}},$$
so pairs $(\widehat{S},\psi)$, where $\widehat{S}\subseteq \widehat{T}$ is a subgroup of finite index and $\psi\in H^2(\widehat{S},\mathbb{C}^{\times})$, are in bijection with finite indecomposable semisimple module categories over $\Rep(T)$, given by $(\widehat{S},\psi)\mapsto \M(\widehat{S},\psi)$. \qed
\end{remark}

\subsubsection{Equivariant structures on $\M(\widehat{S},\Psi)$}\label{sec:Equivariant structures on M}
Fix a $\Rep(T)$-module category $\M(\widehat{S},\Psi)$ as above. For each $k\in K$, consider the $\Rep(T)$-module category $\M(\widehat{S},\Psi)^k$ with action given by
$$\delta_{x}\overline{\ot}^k\delta_{[y]}:=\delta_{x^k}\overline{\ot}\delta_{[y]}=\delta_{[x^ky]},\quad x\in \widehat{T},\,\,[y]\in \widehat{T}/\widehat{S},$$
and module associativity constraint given by
\begin{equation*}
m_{\Psi}^k(\delta_x,\delta_z,\delta_{[y]}):=\Psi(x^k,z^k)([y])\cdot \id
:\delta_{[x^kz^ky]}\xrightarrow{\cong}\delta_{[x^kz^ky]},
\end{equation*}
for every $x,z\in \widehat{T}$ and $[y]\in \widehat{T}/\widehat{S}$. 
Namely, for every $k\in K$, we have
$$\M(\widehat{S},\Psi)^k=\M(\widehat{S},\Psi^k).$$

\begin{theorem}\label{partaunew}
Fix a $\Rep(T)$-module category $\M(\widehat{S},\Psi)$ and a subgroup $L\subseteq K$. 
Then the following hold:
\begin{enumerate}
\item
If $\M(\widehat{S},\Psi)$ admits an $L$-equivariant structure over $\Rep(T)$, then $\widehat{S}$ is $L$-stable.
\item
If $\widehat{S}$ is $L$-stable, then 
$L$-equivariant structures on $\M(\widehat{S},\Psi)$ over $\Rep(T)$ are parametrized by triples $(U,\tau,\mu)$, where 
\begin{equation}\label{Ul0}
U:L\to {\rm Sym}(\widehat{T}/\widehat{S}),\,\,\,\ell\mapsto U_{\ell},
\end{equation}
is a group homomorphism, 
\begin{equation}\label{taufirstt}
\tau: L\to C^1(\widehat{T}\times \widehat{T}/\widehat{S},\mathbb{C}^{\times}),\,\,\,\ell\mapsto \tau_{\ell},
\end{equation}
and
\begin{equation}\label{mufirstt0}
\mu: L\times L\to C^1(\widehat{T}/\widehat{S},\mathbb{C}^{\times}),\,\,\,(k,\ell)\mapsto \mu_{k,\ell},
\end{equation}
are maps, such that for every $i,k,\ell\in L$, $x,z\in \widehat{T}$, and $[y]\in \widehat{T}/\widehat{S}$, the following hold:
\begin{equation}\label{Ul}
U_{\ell}([xy])=x^{\ell}U_{\ell}([y])
\end{equation}
(that is, $U_{\ell}$ is determined by $U_{\ell}([1])\in \widehat{T}/\widehat{S}$),
\begin{equation}\label{taulPsi}
\tau_{\ell}(z,[y])\tau_{\ell}(x,[zy])\Psi(x,z)([y])= \Psi^{\ell}(x,z)(U_{\ell}([y]))\tau_{\ell}(xz,[y]),
\end{equation}
\begin{equation}\label{taul1yS}
\tau_{\ell}(1,[y])=1,
\end{equation}
\begin{equation}\label{Ultaul}
\tau_k(x,U_{\ell}([y]))\tau_{\ell}(x^k,[y])\mu_{k,\ell}([x^{k\ell}y])=\mu_{k,\ell}([y])\tau_{k\ell}(x,[y])\gamma_{k,\ell}(x),
\end{equation}
and 
\begin{equation}\label{mufirstt}
\mu_{i,k}(U_{\ell}([y]))\mu_{ik,\ell}([y])=\mu_{k,\ell}([y])\mu_{i,k\ell}([y]).
\end{equation}
\end{enumerate}
\end{theorem}

\begin{proof}
Suppose that $(U,\tau,\mu)$ is an $L$-equivariant structure on $\M(\widehat{S},\Psi)$ over $\Rep(T)$. Recall \S\ref{sec:Equivariant module categories} that this means the following:

(1) $(U,\tau)=\{(U_{\ell},\tau_{\ell}):\M(\widehat{S},\Psi)\xrightarrow{\simeq} \M(\widehat{S},\Psi^{\ell})\mid \ell\in L\}$ 
is a set of $\Rep(T)$-module equivalences; that is, the following hold: 

(i) For every $\ell\in L$, $U_{\ell}:\widehat{T}/\widehat{S}\xrightarrow{\cong}\widehat{T}/\widehat{S}$ is a set bijection, so that we have an induced set bijection
$$U_{\ell}:\{\delta_{[y]}\mid [y]\in \widehat{T}/\widehat{S}\}\xrightarrow{\cong}\{\delta_{[y]}\mid [y]\in\widehat{T}/\widehat{S}\},\,\,\,\delta_{[y]}\mapsto \delta_{U_{\ell}([y])}.$$

(ii) For every $\ell\in L$, $\tau_{\ell}$ is a natural isomorphism determined by the collection 
$$\{\tau_{\ell}(\delta_x,\delta_{[y]}):\delta_{U_{\ell}([xy])}=U_{\ell}(\delta_{x}\overline{\ot}\delta_{[y]})\xrightarrow{\cong}\delta_{x}\overline{\ot}^{\ell} U_{\ell}(\delta_{[y]})=\delta_{x^{\ell}U_{\ell}([y])}\mid x\in \widehat{T},\,[y]\in \widehat{T}/\widehat{S}\}$$
of nonzero scalars. 
Thus, $\widehat{S}$ must be $L$-stable, and we get (\ref{Ul0})-(\ref{taufirstt}) and (\ref{Ul}). 

Moreover, it follows from (\ref{tauass})-(\ref{tauunit}) that 
\begin{gather*}
(\id_{\delta_{x}}\ot \tau_{\ell}(\delta_z,\delta_{[y]}))\tau_{\ell}(\delta_x,\delta_{z}\overline{\ot}\delta_{[y]})U_{\ell}(m_{\Psi}(\delta_x,\delta_z,\delta_{[y]}))
= m_{\Psi}^{\ell}(\delta_x,\delta_z,\delta_{U_{\ell}(\delta_{[y]})})\tau_{\ell}(\delta_{xz},\delta_{[y]})\\
:U_{\ell}(\delta_{xz}\overline{\ot} \delta_{[y]})\xrightarrow{\cong} \delta_{x}\overline{\ot}^{\ell}(\delta_{z}\overline{\ot}^{\ell} U_{\ell}(\delta_{[y]})),
\end{gather*}
and
\begin{equation*}
l_{U_{\ell}(\delta_{[y]})}\circ\tau_{\ell}(1,\delta_{[y]})=U_{\ell}(l_{[y]}):U_{\ell}(\delta_{[y]})\xrightarrow{\cong} U_{\ell}(\delta_{[y]})
\end{equation*}
for every $x,z\in\widehat{T}$ and $[y]\in \widehat{T}/\widehat{S}$.
Thus, we have that
$$\tau_{\ell}(\delta_z,\delta_{[y]})\tau_{\ell}(\delta_x,\delta_{[zy]})U_{\ell}(m_{\Psi}(\delta_x,\delta_z,\delta_{[y]}))= m_{\Psi}^{\ell}(\delta_x,\delta_z,\delta_{U_{\ell}(\delta_{[y]})})\tau_{\ell}(\delta_{xz},\delta_{[y]}),$$
or equivalently,
$$\tau_{\ell}(\delta_z,\delta_{[y]})\tau_{\ell}(\delta_x,\delta_{[zy]})\Psi(x,z)([y])= \Psi^{\ell}(x,z)(U_{\ell}([y]))\tau_{\ell}(\delta_{xz},\delta_{[y]}),$$
and $\tau_{\ell}(1,\delta_{[y]})=1$, which yield (\ref{taulPsi})-(\ref{taul1yS}) (setting $\tau_{\ell}(x,[y]):=\tau_{\ell}(\delta_x,\delta_{[y]})$).

(2) $\mu=\{\mu_{k,\ell}: (U_{k},\tau_{k})\circ(U_{\ell},\tau_{\ell})^k\xrightarrow{\cong}(U_{k\ell},\tau_{k\ell})\circ(\id_{\M},\Gamma_{k,\ell})\mid k,\ell\in L\}$ is a set of natural isomorphisms of module functors $\M(\widehat{S},\Psi)\to \M(\widehat{S},\Psi^{k\ell})$; that is, 
for every $k,\ell\in L$, we have a collection of nonzero scalars 
$$\{\mu_{k,\ell}(\delta_{[y]}): U_k(U_{\ell}(\delta_{[y]}))=\delta_{U_{k}(U_{\ell}([y]))}\xrightarrow{\cong}\delta_{U_{k\ell}([y])}=U_{k\ell}(\delta_{[y]})\mid [y]\in \widehat{T}/\widehat{S}\},$$
so (\ref{mufirstt0}) holds, 
such that by (\ref{munat0}), we have 
$$
\tau_k(\delta_x,\delta_{U_{\ell}([y])})\tau_{\ell}(\delta_{x^k},\delta_{[y]})\mu_{k,\ell}(x(\overline{\ot}^{\ell})^k \delta_{[y]})=\mu_{k,\ell}(\delta_{[y]})\tau_{k\ell}(\delta_x,\delta_{[y]})\Gamma_{k,\ell}(\delta_x,\delta_{[y]})$$
for every $x\in \widehat{T}$, $[y]\in \widehat{T}/\widehat{S}$ and $k,\ell\in L$,  
and by (\ref{comp00}), we have 
$$
\mu_{i,k}(\delta_{U_{\ell}([y])})\mu_{ik,\ell}(\delta_{[y]})=U_i(\mu_{k,\ell}(\delta_{[y]}))\mu_{i,k\ell}(\delta_{[y]})
$$
for every $[y]\in \widehat{T}/\widehat{S}$ and $i,k,\ell\in L$. 
It follows that  
\begin{equation*}
U_{k\ell}=U_{k}U_{\ell};
\end{equation*}
that is, $U$ is a group homomorphism, and
\begin{equation*}
\tau_k(\delta_x,\delta_{U_{\ell}([y])})\tau_{\ell}(\delta_{x^k},\delta_{[y]})\mu_{k,\ell}(\delta_{[x^{k\ell}y]})=\mu_{k,\ell}(\delta_{[y]})\tau_{k\ell}(\delta_x,\delta_{[y]})\gamma_{k,\ell}(\delta_{x}),
\end{equation*}
so (\ref{Ultaul}) holds (setting $\mu_{k,\ell}([y]):=\mu_{k,\ell}(\delta_{[y]})$ and $\gamma_{k,\ell}(x):=\gamma_{k,\ell}(\delta_{x})$), 
and since $U_i(\mu_{k,\ell}(\delta_{[y]}))=\mu_{k,\ell}(\delta_{[y]})$, we have
\begin{equation*}
\mu_{i,k}(\delta_{U_{\ell}([y])})\mu_{ik,\ell}(\delta_{[y]})=\mu_{k,\ell}(\delta_{[y]})\mu_{i,k\ell}(\delta_{[y]}),
\end{equation*}
so (\ref{mufirstt}) holds too.

Conversely, it is clear that any triple $(U,\tau,\mu)$ satisfying (\ref{taufirstt})-(\ref{mufirstt}) defines an $L$-equivariant structure on $\M(\widehat{S},\Psi)$ over $\Rep(T)$.
\end{proof}

\begin{remark}

(1) (\ref{Ul}) means that $U_{\ell}:\widehat{T}/\widehat{S}\xrightarrow{\cong}(\widehat{T}/\widehat{S})^{\ell}$ is $\widehat{T}$-equivariant.

(2) (\ref{mufirstt}) means that $\mu\in Z^2(L,C^1(\widehat{T}/\widehat{S},\mathbb{C}^{\times}))$, where $L$ acts on $C^1(\widehat{T}/\widehat{S},\mathbb{C}^{\times})$ trivially from the left and via $U$ from the right. \qed
\end{remark}

\subsubsection{Equivariantizations of $\M(\widehat{S},\Psi)$}\label{sec:Equivariantizations of M}
Fix an $L$-equivariant structure $(U,\tau,\mu)$ on $\M(\widehat{S},\Psi)$ over $\Rep(T)$ as in Theorem \ref{partaunew}. Recall \S\ref{sec:Equivariantization of module categories} the equivariantization category $\left(\M(\widehat{S},\Psi),U,\tau,\mu\right)^L$. A nonzero object $(V,v)\in \left(\M(\widehat{S},\Psi),U,\tau,\mu\right)^L$ consists of an object 
\begin{equation}\label{decompV}
V=\bigoplus_{[y]\in {\bf F}}V_{[y]}\ot\delta_{[y]}\in \M(\widehat{S},\Psi);\quad V_{[y]}:={\rm Hom}_{\M(\widehat{S},\Psi)}(\delta_{[y]},V)\ne 0,
\end{equation} 
such that ${\bf F}\subseteq \widehat{T}/\widehat{S}$, $0\ne\dim(V_{[y]})<\infty$ for every $[y]\in {\bf F}$, and   
$$v=\{v_{\ell}: U_{\ell}(V)=\bigoplus_{[y]\in {\bf F}}V_{[y]}\ot\delta_{U_{\ell}([y])}\xrightarrow{\cong} \bigoplus_{[y]\in {\bf F}}V_{[y]}\ot\delta_{[y]}=V\mid {\ell}\in L\}$$ 
is a set of $\M(\widehat{S},\Psi)$-isomorphisms, such that the diagram 
\begin{equation*}
\xymatrix{\bigoplus_{[y]\in {\bf F}}V_{[y]}\ot\delta_{U_k(U_{\ell}([y]))}\ar[d]_{\mu_{k,\ell}(V)}\ar[rr]^{U_k(v_{\ell})}&&\bigoplus_{[y]\in {\bf F}}V_{[y]}\ot\delta_{U_k([y])}\ar[d]^{v_k} &&\\
\bigoplus_{[y]\in {\bf F}}V_{[y]}\ot\delta_{U_{k\ell}([y])}\ar[rr]^{v_{k\ell}}&& \bigoplus_{[y]\in {\bf F}}V_{[y]}\ot\delta_{[y]}&&}
\end{equation*}
commutes for every $k,\ell\in L$ (see (\ref{diagramv})). Thus, the following hold:

(1) The set ${\bf F}$ is $L$-stable (that is, $U_{\ell}({\bf F})={\bf F}$, $\ell\in L$), so for every $\ell\in L$, 
$$V=\bigoplus_{[y]\in {\bf F}}V_{U_{\ell}([y])}\ot\delta_{U_{\ell}([y])},$$   
and $v_{\ell}=\bigoplus_{[y]\in {\bf F}}v_{\ell,[y]}$, where $v_{\ell,[y]}:V_{[y]}\xrightarrow{\cong}V_{U_{\ell}([y])}$ is a $\mathbb{C}$-linear isomorphism. In particular, $\dim(V_{U_{\ell}([y])})=\dim(V_{[y]})$ for every $\ell\in L$ and $[y]\in {\bf F}$.

(2) By naturality of $\mu$, the isomorphism 
$$\mu_{k,\ell}(V): 
\bigoplus_{[y]\in {\bf F}}V_{[y]}\ot\delta_{U_k(U_\ell([y]))}=\bigoplus_{[y]\in {\bf F}}V_{[y]}\ot\delta_{U_{k\ell}([y])}
\xrightarrow{\cong}\bigoplus_{[y]\in {\bf F}}V_{[y]}\ot\delta_{U_{k\ell}([y])}
$$ 
is given by  
$$\mu_{k,\ell}(V)=\bigoplus_{[y]\in {\bf F}}\mu_{k,\ell}([y])\cdot \id_{V_{[y]}}:V_{[y]}\ot\delta_{U_{k\ell}([y])}\xrightarrow{\cong}V_{[y]}\ot\delta_{U_{k\ell}([y])},$$ 
so the diagram
\begin{equation*}
\xymatrix{V_{[y]}\ot\delta_{U_k(U_{\ell}([y]))}\ar[d]_{\mu_{k,\ell}([y])\cdot \id_{V_{[y]}}}\ar[rr]^{U_k(v_{\ell,[y]})}&&V_{[y]}\ot\delta_{U_k([y])}\ar[d]^{v_{k,[y]}} &&\\
V_{[y]}\ot\delta_{U_{k\ell}([y])}\ar[rr]^{v_{k\ell,[y]}}&& V_{[y]}\ot\delta_{[y]}&&}
\end{equation*}
commutes; that is, we have 
\begin{equation}\label{vsat}
v_{k,[y]}\circ U_k(v_{\ell,[y]})=\mu_{k,\ell}(\delta_{[y]})v_{k\ell,[y]}
\end{equation}
in ${\rm Aut}(V_{[y]})$ for all $k,\ell\in L$ and $[y]\in {\bf F}$.
 
\begin{theorem}\label{simpleuiv}
Isomorphism classes of simple objects of $\left(\M(\widehat{S},\Psi),U,\tau,\mu\right)^L$ are parametrized by conjugacy classes of triples $(N,\iota,W)$, where $N\subseteq L$ is a subgroup, $\iota:L/N\xrightarrow{1:1} \widehat{T}/\widehat{S}$ is an $L$-equivariant injective map (that is, for every $f,\ell\in L$, we have $\iota(f[\ell])=U_f(\iota([\ell]))$, so $\iota$ is determined by $\iota([1])\in \widehat{T}/\widehat{S}$), and $W\in {\rm Rep}(N,\mu_N)$ is a simple object, where $\mu_N\in Z^2(N,\mathbb{C}^{\times})$ is defined by $\mu_N(k,\ell)=\mu_{k,\ell}(\delta_{\iota([1])})$.
\end{theorem}

\begin{proof}
Suppose that $(V,v)\in \left(\M(\widehat{S},\Psi),U,\tau,\mu\right)^L$ is a simple object, and consider the decomposition of $V$ as in (\ref{decompV}). Then $L$ acts transitively on ${\bf F}$ via $U$. Let $N\subseteq L$ be the stabilizer of a point in ${\bf F}$, so that we have an $L$-equivariant injective map $\iota:L/N\xrightarrow{\cong}{\bf F}\subseteq  \widehat{T}/\widehat{S}$ (i.e. $\iota(f[\ell])=U_f(\iota([\ell]))$, $f,\ell\in L$). Set $W:=V_{\iota([1])}$. Since $U_k=\id$ for every $k\in N$, it follows from (\ref{vsat}) that  
$$v_{k,\iota([1])}\circ v_{\ell,\iota([1])}=\mu_{k,\ell}(\delta_{\iota([1])})v_{k\ell,\iota([1])}=\mu_N(k,\ell)v_{k\ell,\iota([1])}$$
in ${\rm Aut}(W)$ for every $k,\ell\in N$. Thus, $W\in {\rm Rep}(N,\mu_N)$ is a simple object.

Conversely, given a pair $(N,\iota)$, where $N\subseteq L$ is a subgroup and 
$$\iota:L/N\xrightarrow{1:1} \widehat{T}/\widehat{S}$$ is an $L$-equivariant injective map (where $L$ acts on $\widehat{T}/\widehat{S}$ via $U$), and a simple object $(W,\rho)\in {\rm Rep}(N,\mu_N)$, we assign an object $(V,v)\in \left(\M(\widehat{S},\Psi),U,\tau,\mu\right)^L$ as follows. 
Fix a set 
$$R=\{r_1=1,\dots,r_n\}\subset L$$ of representatives of $L/N$, and define the map  
\begin{equation}\label{defnyi}
\iota:L/N\xrightarrow{1:1} \widehat{T}/\widehat{S},\quad [r_i]\mapsto [y_i].
\end{equation}
For every $\ell\in L$ and $1\le i\le n$, let $r_{\ell(i)}\in R$ and $n_{\ell(i)}\in N$ be the unique elements such that $\ell r_i = r_{\ell(i)}n_{\ell(i)}$. Note that by definition,
$$U_{\ell}([y_i])=U_{\ell}(\iota([r_i]))=\iota([\ell r_i])=\iota([r_{\ell(i)}])=[y_{\ell(i)}].$$ 

Now let
$$V := \bigoplus_{i=1}^n W\ot\delta_{[y_i]}\in \M(\widehat{S},\Psi),$$ 
and let
$$v=\{v_{\ell}:U_{\ell}(V)=\bigoplus_{i=1}^n W\ot\delta_{U_{\ell}([y_i])}\xrightarrow{\cong}\bigoplus_{i=1}^n  W\ot\delta_{[y_i]}=V\mid \ell\in L\}$$
be given by $v_{\ell}=\bigoplus_{i=1}^n v_{\ell,i}$, where 
$$v_{\ell,i}:=\rho(n_{\ell(i)})\ot\id:W\ot\delta_{[y_{\ell(i)}]}\xrightarrow{\cong}W\ot\delta_{[y_{\ell(i)}]}.$$
Then it is straightforward to verify that $(V,v)$ is simple, and that its isomorphism type depends only on the conjugacy class of $(N,\iota)$.

Finally, it is straightforward to verify that the two correspondences are inverse to each other. 
\end{proof}

\begin{remark}\label{setC}
Let ${\rm C}$ denote the set of conjugacy classes of pairs $(N,\iota)$, such that $N\subseteq L$ is a subgroup and $\iota:L/N\xrightarrow{1:1}\widehat{T}/\widehat{S}$ is $L$-equivariant. Then 
Theorem \ref{simpleuiv} can be reformulated by saying that we have an equivalence
$$\left(\M(\widehat{S},\Psi),U,\tau,\mu\right)^L\simeq \bigoplus_{(N,\iota)\in{\rm C}}{\rm Rep}(N,\mu_N)$$
of Abelian categories. \qed
\end{remark}

\begin{corollary}\label{modcatr2r1}
The module category $\left(\M(\widehat{S},\Psi),U,\tau,\mu\right)^L$ has rank $1$ if and only if $N=L$, $\widehat{S}=\widehat{T}$, $U=\id$, and $\mu\in Z^2(L,\mathbb{C}^{\times})$ is nondegenerate. 
\end{corollary}

\begin{proof}
By Theorem \ref{simpleuiv}, $\left(\M(\widehat{S},\Psi),U,\tau,\mu\right)^L$ has rank $1$ if and only if the set ${\rm C}$ from Remark \ref{setC} contains a unique element $(N,\iota)$, which  is the case if and only if $N=L$, $\widehat{S}=\widehat{T}$, $U=\id$, and  $\mu\in Z^2(L,\mathbb{C}^{\times})$ is nondegenerate (see Example \ref{eqvec5}), as claimed.
\end{proof}

Recall \S\ref{sec:Equivariantization of module categories} that $\left(\M(\widehat{S},\Psi),U,\tau,\mu\right)^L$ is a $\Rep(G)=\Rep(T)^K$-module category, and by \cite[Proposition 5.4]{ENO}, every finite indecomposable semisimple module category over $\Rep(G)=\Rep(T)^K$ is of this form.

Let $G_L\subseteq G$ be the subgroup such that $G_L/T=L$ (see (\ref{sesalgr})). Note that $\Rep(G_L)\simeq \Rep(T)^L$, so there is a canonical surjective tensor functor 
$$F_L:\Rep(G)\twoheadrightarrow \Rep(G_L).$$

\begin{corollary}\label{modfactor}
Any finite semisimple module category $\left(\M(\widehat{S},\Psi),U,\tau,\mu\right)^L$ over $\Rep(G)$ factors through $F_L:\Rep(G)\twoheadrightarrow \Rep(G_L)$.
\end{corollary}

\begin{proof}
It is clear from the above (e.g., from Theorem \ref{partaunew}) that $\left(\M(\widehat{S},\Psi),U,\tau,\mu\right)^L$ is a finite indecomposable semisimple module category over $\Rep(G_L)$.
\end{proof}

\subsection{Fiber functors on ${\rm Rep}(G)$}\label{sec:Fiber functors on RepG}
Recall that fiber functors $\Rep(G)\to\Vect$ on $\Rep(G)$ are the same thing as $\Rep(G)$-module categories of rank $1$. 

Fix a module category $\left(\M(\widehat{T},\Psi),\id,\tau,\mu\right)^L$ over $\Rep(G)$ of rank $1$, so that $\mu$ is a nondegenerate $2$-cocycle in $Z^2(L,\mathbb{C}^{\times})$ (see Corollary \ref{modcatr2r1}). Define 
$$\psi\in Z^2(\widehat{T},\mathbb{C}^{\times}),\quad \psi(x,y)=\Psi(x,y,[1]),$$
and set $\V(\psi,L,\tau,\mu):=\left(\M(\widehat{T},\Psi),\id,\tau,\mu\right)^L$. Then
\begin{equation}\label{defnofV}
\V(\psi,L,\tau,\mu)=\Rep(L,\mu)=\Vect
\end{equation}
as Abelian categories. Moreover, if $(W,\rho)\in {\rm Rep}(L,\mu)$ is the unique simple object, then  
the unique simple object ${\rm V}:=(V,v)$ of $\V(\psi,L,\tau,\mu)$ is given by  
\begin{equation*}
V=W\ot_{\mathbb{C}}\delta\in \Vect,
\end{equation*}
where $\delta:=\delta_{[1]}=\mathbb{C}$, and
\begin{equation*}
v=\{v_{\ell}=\rho_{\ell}\ot \id:V=W\ot_{\mathbb{C}}\delta\xrightarrow{\cong} W\ot_{\mathbb{C}}\delta=V\mid \ell\in L\}.
\end{equation*} 
Thus, we will make the identifications $V=W$ and $v=\rho$, and write ${\rm V}:=(W,\rho)$.

Given $\V(\psi,L,\tau,\mu)$, define 
\begin{equation}\label{defbl}
B_L:=\langle\tau(L)\rangle \cap T\subseteq T,
\end{equation} 
where $\langle \tau(L)\rangle\subseteq C^1(\widehat{T},\mathbb{C}^{\times})$ is the subgroup generated by $\tau(L)$ (but not necessarily in $Z^2(L,T)$ since $\tau$ may not take values in $T$).

\begin{proposition}\label{newcofib}
Fix $\V:=\V(\psi,L,\tau,\mu)$ as above. The following hold:
\begin{enumerate} 
\item
$\tau: L\to C^1(\widehat{T},\mathbb{C}^{\times}),\,\,\,\ell\mapsto \tau_{\ell}:=\tau_{\ell}(-,\delta)$,
such that $\tau_{1}=1\in T$, and  
$$\tau_{k}\tau_{\ell}^k=\gamma_{k,\ell}\tau_{k\ell},\quad k,\ell\in L;$$
that is, 
$\gamma_{\mid L}=\partial\tau\in Z^2(L,B_L)$.
\item
For every $\ell\in L$ and $x,y\in \widehat{T}$, $\tau_{\ell}(xy)=\tau_{\ell}(x)\tau_{\ell}(y)(\psi^{-1}\psi^{\ell})(x,y)$; that is, 
$$\psi^{-1}\psi^{\ell}=(\tau_{\ell}\circ {\rm m})(\tau_{\ell}^{-1}\ot \tau_{\ell}^{-1}),\quad \ell\in L.$$ 
In particular, 
$[\psi]\in H^2(\widehat{T},\mathbb{C}^{\times})^L$.
\item
For every $\ell\in L$, $R^{\psi}=(R^{\psi})^{\ell}$. Thus, the support $A_{\psi}\subseteq T$ of $\psi$ is $L$-stable, so $A_{\psi}$ is a normal subgroup of $G_L$. 
\item
The module associativity constraint $m^{\V}$ of $\V$ is determined by 
\begin{equation*}
m^{\V}({\rm X},{\rm Y},{\rm V}):({\rm X}\ot {\rm Y})\ot^{\V} {\rm V}\xrightarrow{\cong}{\rm X}\ot^{\V} ({\rm Y}\ot^{\V} {\rm V})
\end{equation*}
for every ${\rm X}=(X,u),\,{\rm Y}=(Y,w)\in \Rep(G)$; that is, by
$$m_{\psi}(X,Y)\ot \id_{W}:(X\ot_{\mathbb{C}} Y)\ot_{\mathbb{C}} W\xrightarrow{\cong}X\ot_{\mathbb{C}} (Y\ot_{\mathbb{C}} W).$$
\end{enumerate}
\end{proposition}

\begin{proof}
(1) and (2) These follow from Theorem \ref{partaunew}.

(3) This follows from (2).

(4) See \S\ref{sec:Equivariantization of module categories}.
\end{proof}

\begin{proposition}\label{equivclasses}
Let $\V:=\V(\psi,L,\tau,\mu)$ and $\widetilde{\V}:=\V(\widetilde{\psi},\widetilde{L},\widetilde{\tau},\widetilde{\mu})$ be two $\Rep(G)$-module categories of rank $1$, with simple objects ${\rm V}$ and $\widetilde{{\rm V}}$, respectively. Then 
$\V\simeq\widetilde{\V}$ as $\Rep(G)$-module categories if and only if there exists an element $g\in G$ such that
$[\psi]=[\widetilde{\psi}^g]$, $L=g^{-1}\widetilde{L}g$, $\tau=\widetilde{\tau}^g$,\ and $[\mu]=[\widetilde{\mu}^g]$.
\end{proposition}

\begin{proof}
The ``if" part is clear.

For the ``only if" part, suppose that 
$(E,s):\V\xrightarrow{\simeq}\widetilde{\V}$ is a $\Rep(G)$-module equivalence. Then by definition, the following hold:

(1) $E$ is an equivalence of Abelian categories, so for every ${\rm X}=(X,u)\in \Rep(G)$, we have by Proposition \ref{newcofib}(4) that 
$E({\rm X} \ot ^{\V}{\rm V})=F({\rm X})\ot_{\mathbb{C}}\widetilde{{\rm V}}$.

(2) $s$ is a natural isomorphism, determined by 
\begin{equation*}
s_{{\rm X},{\rm V}}:F({\rm X})\ot_{\mathbb{C}}\widetilde{{\rm V}}\xrightarrow{\cong}F({\rm X})\ot_{\mathbb{C}}\widetilde{{\rm V}},\quad {\rm X}\in \Rep(G),
\end{equation*}
such that for every ${\rm X},{\rm Y}\in \Rep(G)$ the following hold:
\begin{equation*}
\begin{split}
& (\id_{{\rm X}}\ot s_{{\rm Y},{\rm V}})\circ s_{{\rm X},{\rm Y} \ot ^{\V}{\rm V}}\circ E(m_{\psi}({\rm X},{\rm Y},{\rm V}))=m_{\widetilde{\psi}}({\rm X},{\rm Y},\widetilde{{\rm V}})\circ s_{{\rm X}\ot {\rm Y},{\rm V}}\\
& :F({\rm X})\ot F({\rm Y})\ot_{\mathbb{C}}\widetilde{{\rm V}} \xrightarrow{\cong} F({\rm X})\ot F({\rm Y})\ot_{\mathbb{C}}\widetilde{{\rm V}},
\end{split}
\end{equation*}
and
\begin{equation*}
l_{\widetilde{{\rm V}}}\circ s_{1,{\rm V}}=E\left(l_{{\rm V}}\right): \widetilde{{\rm V}}\xrightarrow{\cong} 
\widetilde{{\rm V}}.
\end{equation*}
It follows that $s\in {\rm Aut}(F)=(\mathscr{O}(G)^*)^{\times}$ satisfies $(s\ot s)\psi=\widetilde{\psi}\Delta(s)$. Hence, 
$$(s^{-1}\ot s^{-1})\widetilde{\psi}^{-1}(s\ot s)\psi=(s^{-1}\ot s^{-1})\Delta(s)$$ 
is a symmetric Hopf $2$-cocycle for $G$. By \cite{DM}, there exists $\eta\in (\mathscr{O}(T)^*)^{\times}$ such that
$$(s^{-1}\ot s^{-1})\widetilde{\psi}^{-1}(s\ot s)\psi=(\eta^{-1}\ot \eta^{-1})\Delta(\eta).$$
It follows that $g:=\eta^{-1}s\in \mathscr{O}(G)^*$ is a grouplike element, so $g\in G$, and
$$(g\ot g)\psi(g^{-1}\ot g^{-1})=\widetilde{\psi}.$$
The rest of the claim follows now in a straightforward manner.
%This means, that $\psi$ and $\widetilde{\psi}$, viewed as Hopf $2$-cocycles for $G$ via the surjective Hopf algebra map $\mathscr{O}(G)\twoheadrightarrow \mathscr{O}(T)$, are gauge equivalent
\end{proof}

\begin{theorem}\label{modcatr2}
The assignment $(\psi,L,\tau,\mu)\mapsto \V(\psi,L,\tau,\mu)$ determines a one to one correspondence between conjugacy classes of quadruples $(\psi,L,\tau,\mu)$ and equivalence classes of fiber functors $\Rep(G)\to\Vect$. In particular, 
every fiber functor $\Rep(G)\to\Vect$ is classical; that is, the underlying functor is the forgetful functor $F:\Rep(G)\to\Vect$. 
\end{theorem}

\begin{proof}
This follows from Corollary \ref{modcatr2r1} and Proposition \ref{equivclasses}.
\end{proof}

\begin{proposition}\label{finalminimal2}
Fix $\V:=\V(\psi,L,\tau,\mu)$. Let $A_{\psi}\subseteq T$ be the support of $\psi$. Then 
the following hold:
\begin{enumerate}
\item
$A_{\psi}B_L\rtimes_{\eta} L\subseteq G$ is a subgroup.
\item
$\V$ factors through the surjective tensor functor 
$$\Rep(G)\twoheadrightarrow \Rep(A_{\psi}B_L\rtimes_{\eta} L).$$
\end{enumerate} 
\end{proposition}

\begin{proof}
(1) By Proposition \ref{newcofib}(1), we have that $\tau_{\ell}^k=\tau_{k}^{-1}\gamma_{k,\ell}\tau_{k\ell}\in B_L$ for every $k,\ell\in L$, and by Proposition \ref{newcofib}(3), we have that $A_{\psi}$ is $L$-stable. Hence, $A_{\psi}B_L$ is $L$-stable, so $A_{\psi}B_L\rtimes_{\eta} L\subseteq G$ is a subgroup.

(2) It follows from Proposition \ref{newcofib} that $\V$ is a finite indecomposable semisimple module category over $\Rep(A_{\psi}B_L\rtimes_{\eta} L)$.
\end{proof}

\subsection{Hopf $2$-cocycles for $G$} 
Recall \cite{G} that 
gauge equivalence classes of Hopf $2$-cocycles for $T$ are in bijection with elements of the group $H^2(\widehat{T},\mathbb{C}^{\times})$, and gauge equivalence classes of Hopf $2$-cocycles for $K$ are in bijection with conjugacy classes of pairs $(L,\mu)$, where $L\subseteq K$ is a subgroup and $\mu\in Z^2(L,\mathbb{C}^{\times})$ is a nondegenerate $2$-cocycle (see Example \ref{eqvec5}).

By Theorem \ref{red2}(1), Theorem \ref{partaunewgen} applies to the cosemisimple Hopf algebra $H:=\mathscr{O}(T)$. The following result provides more information in this special case.

\begin{theorem}\label{fibredred1}
The following hold:
\begin{enumerate} 
\item
There is a bijection between gauge equivalence classes of Hopf $2$-cocycles for $G$, and conjugacy classes of quadruples $(\psi,L,\tau,\mu)$, such that 
\begin{enumerate}
\item $L\subseteq K$ is a subgroup. 
\item
$\tau: L\to C^1(\widehat{T},\mathbb{C}^{\times}),\,\,\,\ell\mapsto \tau_{\ell}:=\tau_{\ell}(-,\delta)$,
$\tau_{1}=1\in T$, and   
$$\gamma_{\mid L}=\partial\tau\in Z^2(L,B_L).$$ 
\item  
$\psi^{-1}\psi^{\ell}=(\tau_{\ell}\circ {\rm m})(\tau_{\ell}^{-1}\ot \tau_{\ell}^{-1})$ for every $\ell\in L$ (so $[\psi]\in H^2(\widehat{T},\mathbb{C}^{\times})^L$).  
\item 
$\mu\in Z^2(L,\mathbb{C}^{\times})$ is nondegenerate.
\end{enumerate} 
\item
Gauge equivalence classes of {\bf minimal} Hopf $2$-cocycles for $G$ correspond under the above bijection to conjugacy classes of triples $(\psi,\tau,\mu)$, such that the support $A_{\psi}\subseteq T$ of $\psi$ is a normal subgroup of  $G$, $T=A_{\psi}B_K$, and $\mu\in Z^2(K,\mathbb{C}^{\times})$ is nondegenerate (so $K$ is a group of central type).
\end{enumerate}
\end{theorem}

\begin{proof}
(1) This follows from Propositions \ref{newcofib}, \ref{equivclasses}.

(2) If $\V(\psi,L,\tau,\mu)$ is minimal, then by Corollary \ref{modfactor}, $L=K$, and then by Proposition \ref{newcofib}(3), $A_{\psi}$ is normal in $G$, and by Proposition \ref{finalminimal2}, $T=A_{\psi}B_K$.

Conversely, let $\V:=\V(\psi,K,\tau,\mu)$, such that the support $A_{\psi}$ of $\psi$ is normal in $G$, and $T=A_{\psi}B_K$. Suppose that $H\subseteq G$ is a subgroup, such that $\V$ factors through the surjective tensor functor $\Rep(G)\twoheadrightarrow \Rep(H)$ and $\V$ is minimal for $\Rep(H)$. Let $S:=H\cap T$, $\iota:S\hookrightarrow T$, and $\widetilde{L}:=H/S\subseteq G/T=K$. Then it follows from Proposition \ref{newcofib}, applied to $\Rep(H)=\Rep(S)^{\widetilde{L}}$, that we have an equivalence $\V(\psi,K,\tau,\mu)\simeq\V(\widetilde{\psi},\widetilde{L},\widetilde{\tau},\widetilde{\mu})$ of $\Rep(H)$-module categories, such that 
\begin{enumerate}
\item
$\widetilde{\mu}\in Z^2(\widetilde{L},\mathbb{C}^{\times})$ is a nondegenerate $2$-cocycle; 
\item
$\widetilde{\tau}: \widetilde{L}\to C^1(\widehat{S},\mathbb{C}^{\times})$,
such that $\partial\widetilde{\tau}=\gamma_{\mid \widetilde{L}}\in Z^2(\widetilde{L},\langle \widetilde{\tau}(\widetilde{L})\rangle\cap S)$;
\item
$[\widetilde{\psi}]\in H^2(\widehat{S},\mathbb{C}^{\times})^{\widetilde{L}}$, and $\widetilde{\psi}^{-1}\widetilde{\psi}^{m}=(\widetilde{\tau}_{m}\circ {\rm m})(\widetilde{\tau}_{m}^{-1}\ot \widetilde{\tau}_{m}^{-1})$ for every $m\in \widetilde{L}$;
\item
$\widetilde{A}_{\widetilde{\psi}}(\langle \widetilde{\tau}(\widetilde{L})\rangle\cap S)=S$, where $\widetilde{A}_{\widetilde{\psi}}\subseteq S\subseteq T$ is the support of $\widetilde{\psi}$. 
\end{enumerate}
Now since $\V(\psi,K,\tau,\mu)\simeq \V(\iota_*^{\ot 2}\widetilde{\psi},\widetilde{L},\iota_*\widetilde{\tau},\widetilde{\mu})$ as $\Rep(G)$-module categories, it follows from Proposition \ref{equivclasses} that there exists an element $g\in G$ such that
$$[\psi]=[(\iota_*^{\ot 2}\widetilde{\psi})^{g}],\quad K=g^{-1}\widetilde{L}g,\quad \tau=(\iota_*\widetilde{\tau})^g,\quad [\mu]=[\widetilde{\mu}^g].$$ 
In particular, $\widetilde{L}=K$, so we have $A_{\psi}=g^{-1}A_{\psi}g=\widetilde{A}_{\widetilde{\psi}}$. Thus, $A_{\psi}(\langle \widetilde{\tau}(K)\rangle\cap S)=S$. But, $\widetilde{\tau}(K)=\tau(K)$, hence $S=A_{\psi}(\langle \widetilde{\tau}(K)\rangle\cap S)=A_{\psi}(\langle \tau(K)\rangle\cap T)=A_{\psi}B_K=T$, so $H=G$.
\end{proof}

\begin{example}\label{partau1}
Each quadruple $(\psi,1,1,1)$ corresponds to the lifting of $\psi$ to $G$ via the surjective Hopf algebra map $\O(G)\twoheadrightarrow \O(T)$, or equivalently, the lifting of the fiber functor ${\rm Rep}(T)\xrightarrow{(F,\psi)} \Vect$ to ${\rm Rep}(G)$ via the surjective tensor functor ${\rm Rep}(G)\twoheadrightarrow {\rm Rep}(T)$. \qed
\end{example}

\begin{corollary}\label{directproductcase}
Let $G=T\times K$ be a direct product. Then the following hold:
\begin{enumerate}
\item
Gauge equivalence classes of Hopf $2$-cocycles for $T\times K$ are in bijection with conjugacy classes of quadruples $(\psi,L,\tau,\mu)$, such that $\psi\in Z^2(\widehat{T},\mathbb{C}^{\times})$, $L\subseteq K$ is a subgroup, $\tau:L/[L,L]\to T$ is a group homomorphism, and $\mu\in Z^2(L,\mathbb{C}^{\times})$ is nondegenerate.
\item
Gauge equivalence classes of {\bf minimal} Hopf $2$-cocycles for $T\times K$ are in bijection with conjugacy classes of triples $(\psi,\tau,\mu)$, such that $\psi\in Z^2(\widehat{T},\mathbb{C}^{\times})$ (with support $A_{\psi}$), $\tau:K/[K,K]\to T$ is a group homomorphism, $T$ is generated by $A_{\psi}$ and ${\rm Im}(\tau)$, and $\mu\in Z^2(K,\mathbb{C}^{\times})$ is nondegenerate. \qed
\end{enumerate}
\end{corollary}

\begin{proof}
Since in this case $\gamma=1$, and all conjugation actions are trivial, we have that $\tau_{\ell}^k=\tau_{\ell}$ and $\psi^{\ell}=\psi$ for every $k,\ell\in L$. Hence, $\tau:L\to T$ is a group homomorphism. Thus, the statement follows from Theorem \ref{fibredred1}.
\end{proof}

\begin{example}(Commutative direct products) \label{directprodAbuctcase}
Let $G=T\times K$ be {\bf commutative}. Let $Q$ denote the set of conjugacy classes $[\psi,L,\tau,\mu]$ of quadruples $(\psi,L,\tau,\mu)$, where $L\subseteq K$ is a subgroup, $\tau: L\to T$ is a group homomorphism, $\psi\in Z^2(\widehat{T},\mathbb{C}^{\times})$, and $\mu\in Z^2(L,\mathbb{C}^{\times})$ is nondegenerate. Note that the class $[\psi,L,\tau,\mu]\in Q$ consists of all quadruples $(\psi',L,\tau,\mu')$ such that $[\psi]=[\psi']$ and $[\mu]=[\mu']$. 
Then by Corollary \ref{directproductcase}, the following hold:
\begin{enumerate}
\item
Gauge equivalence classes of Hopf $2$-cocycles for $T\times K$ are in bijection with conjugacy classes of quadruples $(\psi,L,\tau,\mu)$, such that $\psi\in Z^2(\widehat{T},\mathbb{C}^{\times})$, $L$ is a subgroup of $K$, $\tau:L\to T$ is a group homomorphism, and $\mu\in Z^2(L,\mathbb{C}^{\times})$ is nondegenerate.
\item
Gauge equivalence classes of {\bf minimal} Hopf $2$-cocycles for $T\times K$ are in bijection with conjugacy classes of triples $(\psi,\tau,\mu)$, such that $\psi\in Z^2(\widehat{T},\mathbb{C}^{\times})$ (with support $A_{\psi}$), $\tau:K\to T$ is a group homomorphism, $T=A_{\psi}{\rm Im}(\tau)$, and $\mu\in Z^2(K,\mathbb{C}^{\times})$ is nondegenerate.
\end{enumerate}

On the other hand, recall \cite{G2} that Hopf $2$-cocycles 
$$J:\O(T\times K)^{\ot 2}=\mathbb{C}[(\widehat{T\times K})^{2}]\to \mathbb{C}$$
for $T\times K$ are nothing but $2$-cocycles $J\in Z^2(\widehat{T\times K},\mathbb{C}^{\times})$, so 
the group of gauge equivalence classes of Hopf $2$-cocycles for $T\times K$
is the Abelian group $H^2(\widehat{T\times K},\mathbb{C}^{\times})$. Recall that we have a group isomorphism  
\begin{equation}\label{twaltbil}
\mathscr{R}:H^2(\widehat{T\times K},\mathbb{C}^{\times})\xrightarrow{\cong}\Hom(\wedge^2(\widehat{T\times K}),\mathbb{C}^{\times}),\,\,\,[J]\mapsto R^J:=J_{21}^{-1}J,
\end{equation}
and {\bf minimal} Hopf $2$-cocycles for $T\times K$ correspond under (\ref{twaltbil}) to {\bf nondegenerate} alternating bicharacters on $\widehat{T\times K}$.

Thus, by (1) and (\ref{twaltbil}), there are bijections
\begin{equation}\label{onehand}
H^2(\widehat{T\times K},\mathbb{C}^{\times})\cong Q,
\end{equation}
and by (2), $[\psi,L,\tau,\mu]\in Q$ corresponds to a minimal Hopf $2$-cocycle for $T\times K$ if and only if $L=K$, and $T=A_{\psi}{\rm Im}(\tau)$. 

For the sake of completeness and reader's convenience, we provide in the appendix below explicit bijections ${\rm Q}:H^2(\widehat{T\times K},\mathbb{C}^{\times})\to Q$ and ${\rm J}:Q\to H^2(\widehat{T\times K},\mathbb{C}^{\times})$ which are inverse to each other. \qed
\end{example}

\section{Appendix: $H^2(\widehat{T\times K},\mathbb{C}^{\times})\cong Q$}\label{sec:appendix}
Retain the setting and notation from Example \ref{directprodAbuctcase}.

\subsection{The map ${\rm Q}:H^2(\widehat{T\times K},\mathbb{C}^{\times})\to Q$} 
Let $n\ge 1$ be the rank of $T$. Fix a basis $x_1,\dots,x_n$ for $\widehat{T}$, so that the map
$$\mathbb{Z}^n\xrightarrow{\cong}\widehat{T},\,\,\,(q_1,\dots, q_n)\mapsto x_1^{q_1}\cdots x_n^{q_n},$$
is a group isomorphism. The following result is well known.

\begin{lemma}\label{variousgrisoms}
The following hold:
\begin{enumerate}
\item
We have a group isomorphism
$$\Phi:\wedge^2(\widehat{T}\times \widehat{K})\xrightarrow{\cong}\wedge^2\widehat{T}\times \widehat{K}^n\times \wedge^2\widehat{K},$$
$$(x,\alpha)\wedge (y,\beta)\mapsto (x\wedge y,(\beta^{q_1}\alpha^{-r_1},\dots,\beta^{q_n}\alpha^{-r_n}),\alpha\wedge\beta),$$
for every $x=x_1^{q_1}\cdots x_n^{q_n},\,y=x_1^{r_1}\cdots x_n^{r_n}\in \widehat{T}$, and 
$\alpha,\beta\in \widehat{K}$.
\item
The inverse group isomorphism 
$$\Phi^{-1}:\wedge^2\widehat{T}\times \widehat{K}^n\times \wedge^2\widehat{K}\xrightarrow{\cong}\wedge^2(\widehat{T}\times \widehat{K})$$
is given for every $x,y\in \widehat{T}$, and $\alpha_1,\dots,\alpha_n,\alpha,\beta\in \widehat{K}$ 
by 
\begin{equation*}
\begin{split}
&\left(x\wedge y,(\alpha_1,\dots,\alpha_n),\alpha\wedge\beta\right)\mapsto \\
& (x,1)\wedge (y,1)+(x_1,1)\wedge (1,\alpha_1)+\cdots +(x_n,1)\wedge (1,\alpha_n)+(1,\alpha)\wedge(1,\beta).
\end{split}
\end{equation*}
\item
We have a group isomorphism
\begin{equation*}
\begin{split}
& \Phi^{-1}_*\mathscr{R}:H^2(\widehat{T}\times \widehat{K},\mathbb{C}^{\times})\xrightarrow{\cong} \Hom(\wedge^2\widehat{T},\mathbb{C}^{\times})\times \Hom(\widehat{K}^n,\mathbb{C}^{\times})\times \Hom(\wedge^2\widehat{K},\mathbb{C}^{\times}),\\
& [J]\mapsto \left(R^J_{\mid \wedge^2(\widehat{T}\times 1)},\left(R^J\left((x_1,1),(1,-)\right),\dots,R^J\left((x_n,1),(1,-)\right)\right),R^J_{\mid \wedge^2(1\times \widehat{K})}\right).
\end{split}
\end{equation*}
\item
For every $x=x_1^{q_1}\cdots x_n^{q_n},\,y=x_1^{r_1}\cdots x_n^{r_n}\in \widehat{T}$ and 
$\alpha,\beta\in \widehat{K}$, we have that
$$\Phi_*:\Hom(\wedge^2\widehat{T},\mathbb{C}^{\times})\times \Hom(\widehat{K}^n,\mathbb{C}^{\times})\times \Hom(\wedge^2\widehat{K},\mathbb{C}^{\times})\xrightarrow{\cong}\Hom(\wedge^2(\widehat{T}\times \widehat{K}),\mathbb{C}^{\times})$$
is given by
\begin{equation*}
\begin{split}
& \Phi_*\left(B_1,(k_1,\dots,k_n),B_2\right)((x,\alpha),(y,\beta))\\
&=B_1(x,y) \langle ((q_1,\dots, q_n)\ot\beta)((r_1,\dots, r_n)\ot\alpha^{-1}),(k_1,\dots,k_n)\rangle B_2(\alpha,\beta)\\
&=B_1(x,y) \beta(k_1^{q_1}\cdots k_n^{q_n})\alpha^{-1}(k_1^{r_1}\cdots k_n^{r_n}) B_2(\alpha,\beta).
\end{split}
\end{equation*}
\end{enumerate}
\end{lemma}

\begin{proof}
(1) The isomorphism $\Phi$ is given by the following composition:
$$\wedge^2(\widehat{T}\times \widehat{K})\xrightarrow{\cong}\wedge^2\widehat{T}\times (\widehat{T}\ot_{\mathbb{Z}} \widehat{K})\times \wedge^2\widehat{K}\xrightarrow{\cong}\wedge^2\widehat{T}\times \widehat{K}^n\times \wedge^2\widehat{K},$$
$$(x,\alpha)\wedge (y,\beta)\mapsto (x\wedge y,x\ot\beta -y\ot\alpha,\alpha\wedge\beta)\mapsto (x\wedge y,(\beta^{q_1}\alpha^{-r_1},\dots,\beta^{q_n}\alpha^{-r_n}),\alpha\wedge\beta),$$
for every $x=x_1^{q_1}\cdots x_n^{q_n},\,y=x_1^{r_1}\cdots x_n^{r_n}\in \widehat{T}$, and 
$\alpha,\beta\in \widehat{K}$.

(2) The inverse isomorphism $\Phi^{-1}$ is given by the following composition:
\begin{equation*}
\begin{split}
& \wedge^2\widehat{T}\times \widehat{K}^n\times \wedge^2\widehat{K}\xrightarrow{\cong}\wedge^2\widehat{T}\times (\widehat{T}\ot_{\mathbb{Z}} \widehat{K})\times \wedge^2\widehat{K}\xrightarrow{\cong}\wedge^2(\widehat{T}\times \widehat{K}),\\
&(x\wedge y,(\alpha_1,\dots,\alpha_n),\alpha\wedge\beta)\mapsto (x\wedge y,x_1\ot \alpha_1+\cdots +x_n\ot \alpha_n,\alpha\wedge\beta)\\
&\mapsto (x,1)\wedge (y,1)+(x_1,1)\wedge (1,\alpha_1)+\cdots +(x_n,1)\wedge (1,\alpha_n)+(1,\alpha)\wedge(1,\beta).
\end{split}
\end{equation*}

(3) This is straightforward using (\ref{twaltbil}) and (2).

(4) This is straightforward.
\end{proof}

Now fix a class $[J]\in H^2(\widehat{T}\times\widehat{K},\mathbb{C}^{\times})$ of Hopf $2$-cocycles for $T\times K$. 
Set 
\begin{equation}\label{theelementskfromJ0}
[\psi]=[J_{\mid \wedge^2(\widehat{T}\times 1)}]\in H^2(\widehat{T},\mathbb{C}^{\times}),\quad [\xi]=J_{\mid \wedge^2(1\times \widehat{K})}\in H^2(\widehat{K},\mathbb{C}^{\times}),
\end{equation} 
\begin{equation}\label{theelementskfromJ}
k_1:=R^J\left((x_1,1),(1,-)\right),\dots,k_n:=R^J\left((x_n,1),(1,-)\right)\in \widehat{\widehat{K}}=K.
\end{equation}
Then by (\ref{twaltbil}), $R^J_{\mid \wedge^2(\widehat{T}\times 1)}=R^{\psi}$ and $R^J_{\mid \wedge^2(1\times \widehat{K})}=R^{\xi}$,
and by Lemma \ref{variousgrisoms}(3), 
$$(\Phi^{-1}_*\mathscr{R})([J])=(R^{\psi},(k_1,\dots,k_n),R^{\xi}).$$  
In particular, we have a group homomorphism $\phi:\widehat{K}\to T$, determined by  
\begin{equation}\label{theelementtheta}
\widehat{\phi}:\widehat{T}\to \widehat{\widehat{K}}=K,\,\,\,x_i\mapsto k_i,
\end{equation}
where the $k_i$s are defined in (\ref{theelementskfromJ}). 

Next let $L\subseteq K$ be the subgroup, such that $L^{\perp}\subseteq \widehat{K}$ is the support of $\xi$; that is,
\begin{equation}\label{thesubgroupL}
{\rm Rad}(R^{\xi})=\widehat{K/L}=L^{\perp}\subseteq \widehat{K}.
\end{equation}
Fix a group isomorphism $D:K\xrightarrow{\cong}\widehat{K}$, and set 
\begin{equation}\label{theelementmu}
\mu:=\xi\circ D^2\in Z^2(K,\mathbb{C}^{\times}). 
\end{equation}
Since $\mu$ is supported on $L$, we may assume that $\mu\in Z^2(L,\mathbb{C}^{\times})$ is nondegenerate. 
Also, use (\ref{theelementtheta}) to define the group homomorphism 
\begin{equation}\label{theelementtau}
\tau:=\phi\circ D\circ \iota:L\to T,
\end{equation}
where $\iota:L\hookrightarrow K$ is the inclusion morphism. 

We have thus obtained a well defined map
\begin{equation}\label{themapQ}
{\rm Q}:H^2(\widehat{T}\times \widehat{K},\mathbb{C}^{\times})\to Q,\quad [J]\mapsto [\psi,L,\tau,\mu],
\end{equation}
where $\psi$, $L$, $\mu$ and $\tau$ are defined in (\ref{theelementskfromJ0}), (\ref{thesubgroupL}), (\ref{theelementmu}) and (\ref{theelementtau}), respectively. 

\subsection{The inverse map ${\rm J}:Q\to H^2(\widehat{T\times K},\mathbb{C}^{\times})$}
Fix a class $[\psi,L,\tau,\mu]\in Q$. Let $\iota:L\hookrightarrow K$ be the inclusion morphism, and fix a group isomorphism $D_L:\widehat{L}\xrightarrow{\cong}L$. Then the $2$-cocycle
\begin{equation}\label{theelementxi}
\xi:=\mu\circ (D_L\circ\widehat{\iota})^2\in Z^2(\widehat{K},\mathbb{C}^{\times})
\end{equation} 
is supported on $L^{\perp}\subseteq \widehat{K}$, and we can use $\tau: L\to T$ to define the elements 
\begin{equation}\label{theelementsk}
k_1:=D_L\widehat{\tau}(x_1),\dots,k_n:=D_L\widehat{\tau}(x_n)\in L\subseteq K.
\end{equation}
Set ${\bf k}:=(k_1,\dots,k_n)\in K^n$.

Now recall from the theory of Schur multipliers that we have group isomorphisms
\begin{equation*}
\begin{split}
& H^2(\widehat{T},\mathbb{C}^{\times})\times K^n\times H^2(\widehat{K},\mathbb{C}^{\times})\xrightarrow{\cong} H^2(\widehat{T},\mathbb{C}^{\times})\times \widehat{K}^n\times H^2(\widehat{K},\mathbb{C}^{\times})\\
& \xrightarrow{\cong}H^2(\widehat{T},\mathbb{C}^{\times})\times (\widehat{T}\ot_{\mathbb{Z}} \widehat{K})\times H^2(\widehat{K},\mathbb{C}^{\times})\xrightarrow{\cong} H^2(\widehat{T}\times \widehat{K},\mathbb{C}^{\times}),
\end{split}
\end{equation*}
such that the composition isomorphism 
\begin{equation}\label{otherhand}
\Theta:H^2(\widehat{T},\mathbb{C}^{\times})\times K^n\times H^2(\widehat{K},\mathbb{C}^{\times})
\xrightarrow{\cong}H^2(\widehat{T}\times \widehat{K},\mathbb{C}^{\times})
\end{equation}
is given by $\Theta([\psi],{\bf k},[\xi])=[\theta([\psi],{\bf k},[\xi])]$, where 
\begin{equation}\label{explicitTheta}
\theta([\psi],{\bf k},[\xi])((x,\alpha),(y,\beta))=\psi(x,y)\beta(k_1^{q_1}\cdots k_n^{q_n})\alpha^{-1}(k_1^{r_1}\cdots k_n^{r_n})\xi(\alpha,\beta)
\end{equation}
for every $x=x_1^{q_1}\cdots x_n^{q_n},\,y=x_1^{r_1}\cdots x_n^{r_n}\in \widehat{T}$ and $\alpha,\beta\in \widehat{K}$.

We have thus obtained a well defined map
\begin{equation}\label{finalhand}
{\rm J}:Q\to H^2(\widehat{T}\times \widehat{K},\mathbb{C}^{\times}),\,\,\,[\psi,L,\tau,\mu]\mapsto \Theta([\psi],{\bf k},[\xi]),
\end{equation}
where $\xi$ and ${\bf k}=(k_1,\dots,k_n)$ are defined in (\ref{theelementxi}) and (\ref{theelementsk}), respectively.

\begin{proposition} 
The maps ${\rm Q}$ (\ref{onehand}) and ${\rm J}$ (\ref{finalhand}) are inverse to each other.
\end{proposition}

\begin{proof}
We first verify that ${\rm J}({\rm Q}([J]))=[J]$, or equivalently by Lemma \ref{variousgrisoms}(3), that
$$(\Phi^{-1}_*\mathscr{R})({\rm J}({\rm Q}([J])))=(\Phi^{-1}_*\mathscr{R})([J]).$$
Indeed, let $[J]\in H^2(\widehat{T}\times \widehat{K},\mathbb{C}^{\times})$ with {\bf normal} $J$, and let ${\rm J}({\rm Q}([J]))=[I]$ with {\bf normal} $I$. Then by definition, we have
\begin{equation*}
\begin{split}
& (\Phi^{-1}_*\mathscr{R})({\rm J}({\rm Q}([J])))=(\Phi^{-1}_*\mathscr{R})([I])\\
&=\left(R^{I}_{\mid \wedge^2(\widehat{T}\times 1)},\left(R^{I}\left(\left(x_1,1\right),\left(1,-\right)\right),\dots,R^{I}\left(\left(x_n,1\right),\left(1,-\right)\right)\right),R^{I}_{\mid \wedge^2(1\times \widehat{K})}\right)\\
&=\left(R^{I}_{\mid \wedge^2(\widehat{T}\times 1)},\left(I\left(\left(x_1,1\right),\left(1,-\right)\right),\dots,I\left(\left(x_n,1\right),\left(1,-\right)\right)\right),R^{I}_{\mid \wedge^2(1\times \widehat{K})}\right).
\end{split}
\end{equation*}
But if ${\rm Q}([J])=[\psi,L,\tau,\mu]$, 
then for every $(x,\alpha),(y,\beta)\in \widehat{T}\times \widehat{K}$, we have by \ref{explicitTheta})-(\ref{finalhand}) that
\begin{equation*}
\begin{split}
& R^{I}((x,1),(y,1))=R^{{\rm J}([\psi,L,\tau,\mu])}((x,1),(y,1))=R^{\theta([\psi],{\bf k},[\xi])}((x,1),(y,1))\\
& =\frac{\theta([\psi],{\bf k},[\xi])((x,1),(y,1))}{\theta([\psi],{\bf k},[\xi])((y,1),(x,1))}=\frac{\psi(x,y)\xi(1,1)}{\psi(y,x)\xi(1,1)}=R^{\psi}(x,y)=R^{J}((x,1),(y,1)),
\end{split}
\end{equation*}
\begin{equation*}
\begin{split}
& R^{I}((x_i,1),(1,\beta))=R^{{\rm J}([\psi,L,\tau,\mu])}((x_i,1),(1,\beta))=R^{\theta([\psi],{\bf k},[\xi])}((x_i,1),(1,\beta))\\
& =\frac{\theta([\psi],{\bf k},[\xi])((x_i,1),(1,\beta))}{\theta([\psi],{\bf k},[\xi])((1,\beta),(x_i,1))}=\psi(x_i,1)\beta(k_i)\xi(1,\beta)=\beta(k_i)=
\beta(D_L\widehat{\tau}(x_i))\\
&=J((x_i,1),(1,\beta))=R^J((x_i,1),(1,\beta))
\end{split}
\end{equation*}
(here we assume that $\theta([\psi],{\bf k},[\xi])$ is normal), and
\begin{equation*}
\begin{split}
& R^{I}((1,\alpha),(1,\beta))=R^{{\rm J}[\psi,L,\tau,\mu]}((1,\alpha),(1,\beta))=R^{\theta([\psi],{\bf k},[\xi])}((1,\alpha),(1,\beta))\\
& =\frac{\theta([\psi],{\bf k},[\xi])((1,\alpha),(1,\beta))}{\theta([\psi],{\bf k},[\xi])((1,\beta),(1,\alpha))}=\frac{\psi(1,1)\xi(\alpha,\beta)}{\psi(1,1)\xi(\beta,\alpha)}=R^{\xi}(\alpha,\beta)=R^{J}((1,\alpha),(1,\beta)),
\end{split}
\end{equation*}
so the claim follows.

Finally, the fact that ${\rm Q}({\rm J}([\psi,L,\tau,\mu]))=[\psi,L,\tau,\mu]$ is verified similarly.
\end{proof}

\subsection{Minimal Hopf $2$-cocycles} 
%We now consider minimal Hopf $2$-cocycles for $G=T\times K$.
Fix a class $[\psi,L,\tau,\mu]\in Q$. Let
\begin{equation}\label{ndB}
B:=R^{{\rm J}([\psi,L,\tau,\mu])}=R^{\theta([\psi],\mathbf{k},[\xi])}\in\Hom(\wedge^2(\widehat{T}\times \widehat{K}),\mathbb{C}^{\times});
\end{equation}
that is, $B=\mathscr{R}({\rm J}([\psi,L,\tau,\mu]))$. 

\begin{proposition}
The following hold:
\begin{enumerate}
\item
For every $(x,\alpha),(y,\beta)\in \widehat{T}\times \widehat{K}$, we have
\begin{equation*}
B((x,\alpha),(y,\beta))=R^{\psi}(x,y)\beta(D_L\widehat{\tau}(x))\alpha^{-1}(D_L\widehat{\tau}(y))R^{\xi}(\alpha,\beta).
\end{equation*}
\item
${\rm Rad}(B)=\{(x,\alpha)\in \widehat{T}\times \widehat{K}\mid R^{\psi}(x,-)=\alpha D_L\widehat{\tau}\,\,\,\&\,\,\,R^{\xi}(\alpha,-)=D_L\widehat{\tau}(x)^{-1}\}$.
\item  
Let $M:=D_L({\rm Im}(\widehat{\tau}))\subseteq L$. Then $B$ is nondegenerate if and only if 
$${\rm Ker}(\widehat{\tau})\cap {\rm Rad}(R^{\psi})=1\text{ and }M^{\perp}\cap {\rm Rad}(R^{\xi})=\widehat{K/M}\cap \widehat{K/L}=1,$$
if and only if,  
$$L=K\text{ and }\tau(K)^{\perp}\cap {\rm Rad}(R^{\psi})=1.$$
(Note that if $L=K$, then ${\rm Rad}(R^{\xi})=1$ is automatic.)
\end{enumerate}
\end{proposition}

\begin{proof}
(1) By Lemma \ref{variousgrisoms}(4), for every $(x,\alpha),(y,\beta)\in \widehat{T}\times \widehat{K}$, we have
\begin{equation*}
\begin{split}
& B((x,\alpha),(y,\beta))\\
& =R^{\psi}(x,y)\beta(D_L\widehat{\tau}(x_1)^{q_1}\cdots D_L\widehat{\tau}(x_n)^{q_n})\alpha^{-1}(D_L\widehat{\tau}(x_1)^{r_1}\cdots D_L\widehat{\tau}(x_n)^{r_n}) R^{\xi}(\alpha,\beta)\\
& =R^{\Psi}(x,y) \beta(D_L\widehat{\tau}(x))\alpha^{-1}(D_L\widehat{\tau}(y)) R^{\xi}(\alpha,\beta),
\end{split}
\end{equation*}
as claimed.

(2) Suppose that $(x,\alpha)\in {\rm Rad}(B)$; 
that is, 
$$B((x,\alpha),(y,\beta))=1$$
for every $(y,\beta)\in \widehat{T}\times \widehat{K}$.  
Then by (1), for $\beta=1$, we have that
$$R^{\psi}(x,y)=\alpha(D_L\widehat{\tau}(y))$$
for every $y\in \widehat{T}$; that is, 
$$R^{\psi}(x,-)=\alpha D_L\widehat{\tau}\in \Hom(\widehat{T},\mathbb{C}^{\times})=T.$$
Also by (1), for $y=1$, we have that
$$R^{\xi}(\alpha,\beta)=\beta^{-1}(D_L\widehat{\tau}(x))$$
for every $\beta\in \widehat{K}$; that is, 
$$R^{\xi}(\alpha,-)=D_L\widehat{\tau}(x)^{-1}\in L\subseteq K=\Hom(\widehat{K},\mathbb{C}^{\times}).$$

The converse is immediate from (1). 

(3) Assume that $B$ is nondegenerate. Let $1\ne x\in \widehat{T}$. Since $(x,1)\notin {\rm Rad}(B)$, it follows from (2) that either $R^{\psi}(x,-)\ne 1$, 
or  
$\widehat{\tau}(x)\ne 1$. It follows that for every $1\ne x\in {\rm Ker}(\widehat{\tau})$, we have $R^{\psi}(x,-)\ne 1$, and for every $1\ne x\in {\rm Rad}(R^{\psi})$, we have $\widehat{\tau}(x)\ne 1$. Thus, ${\rm Ker}(\widehat{\tau})\cap {\rm Rad}(R^{\psi})=1$.

Similarly, let $1\ne \alpha\in \widehat{K}$. Since $(1,\alpha)\notin {\rm Rad}(B)$, it follows from (2) that either $\alpha D_L\widehat{\tau}\ne 1$,  
or $R^{\xi}(\alpha,-)\ne 1$. Hence, for every $1\ne \alpha\in \widehat{K}$ with $\alpha D_L\widehat{\tau}=1$, we have $R^{\xi}(\alpha,-)\ne 1$, 
and for every $1\ne \alpha\in {\rm Rad}(R^{\xi})$, we have $\alpha D_L\widehat{\tau}\ne 1$. Thus, $M^{\perp}\cap {\rm Rad}(R^{\xi})=1$.

Conversely, assume that ${\rm Ker}(\widehat{\tau})\cap {\rm Rad}(R^{\Psi})=1$ and $M^{\perp}\cap {\rm Rad}(R^{\xi})=1$. Let $(x,\alpha)\in {\rm Rad}(B)$. Then since $\alpha$ has a finite order, say $n$, $(x^n,1)\in {\rm Rad}(B)$, too.
But then by (2), we have that
$$R^{\Psi}(x^n,-)=1\,\,\,\text{and}\,\,\,\widehat{\tau}(x^n)=1;$$
that is, $x^n\in {\rm Ker}(\widehat{\tau})\cap {\rm Rad}(R^{\Psi})$. Hence, $x^n=1$, so $x=1$.

Thus, $(1,\alpha)\in {\rm Rad}(B)$, so by (2), we have
$$1=\alpha D_L\widehat{\tau}\,\,\,\text{and}\,\,\,R^{\xi}(\alpha,-)=1;$$
that is, $\alpha\in M^{\perp}\cap {\rm Rad}(R^{\xi})$. Hence, $\alpha=1$, too.
\end{proof}

\begin{corollary}
Let $A\subseteq T$ be the support of $\psi$, so that $\widehat{A}=\widehat{T}/{\rm Rad}(R^{\psi})$. Then $B$ (\ref{ndB}) is nondegenerate if and only if $L=K$ and $A\tau(K)=T$.
\end{corollary}

\begin{proof}
Assume that $B$ is nondegenerate. Then $L=K$. 
Thus, we have an injective group homomorphism 
$$\widehat{T/\tau(K)}=\tau(K)^{\perp}\xrightarrow{1:1}\widehat{T}/{\rm Rad}(R^{\Psi})=\widehat{A},$$
or equivalently, a surjective group homomorphism 
$$A \twoheadrightarrow T/\tau(K),$$
which implies that $A\tau(K)=T$.

The proof of the converse is clear.
\end{proof}

\begin{thebibliography}{ABCD}

\bibitem
[A]{A} E. Abe. Hopf algebras. Translated from the Japanese by Hisae
Kinoshita and Hiroko Tanaka. Cambridge Tracts in Mathematics. 
{\bf 74}, Cambridge University Press, Cambridge-New York, 1980. xii+284 pp.

\bibitem
[DM]{DM} P. Deligne and J. Milne. Tannakian Categories. {\em Lecture Notes in Mathematics} {\bf 900}
(1982), 101--228.

\bibitem
[EG1]{EG1} P. Etingof and S. Gelaki. On cotriangular Hopf algebras, {\em American Journal of Mathematics} {\bf 123} (2001), 699--713.

\bibitem
[EG2]{EG2} P. Etingof and S. Gelaki. Quasisymmetric and unipotent tensor categories, {\em Math. Res. Lett.} {\bf 15} (2008), no. 5, 857--866.

\bibitem
[ENO]{ENO} P. Etingof, D. Nikshych, and V. Ostrik. Weakly group-theoretical and solvable fusion categories. {\em Advances in Mathematics} {\bf 226} (2011), 
176--205.

\bibitem 
[EGNO]{EGNO} 
P. Etingof, S. Gelaki, D. Nikshych, and V. Ostrik. Tensor categories. {\em AMS Mathematical Surveys and Monographs book series} {\bf 205} (2015), 362 pp. 

\bibitem
[G1]{G} S. Gelaki. Twisting of affine algebraic groups, I. {\em International Mathematics Research Notices} 
{\bf 16} (2015), 7552--7574.

\bibitem
[G2]{G2} S. Gelaki. Module categories over affine group schemes. {\em Quantum Topology} (2015), no. 1, 1--37.

\bibitem
[Mon]{M} S. Montgomery. Hopf algebras and their actions on rings. {\em CBMS Regional Conference Series in Mathematics} {\bf 82} (1993), 238 pp.

\bibitem
[Mov]{Mo} 
M. V. Movshev, Twisting in group algebras of finite groups. (Russian) {\em Funktsional. Anal. i Prilozhen.} {\bf 27} (1993), no. 4, 17--23, 95; translation in {\em Funct. Anal. Appl.} {\bf 27} (1993), no. 4, 240--244 (1994). 

\bibitem
[S]{S} T.A. Springer. Linear algebraic groups (2nd ed.). Progress in Mathematics {\bf 9} (1998).
\end{thebibliography}

\end{document}


\begin{example}\label{trivialgamma}
$\gamma$ is trivial. \qed
\end{example}

\begin{example}\label{torus}
$T$ is a torus. \qed
\end{example}

\begin{example}
Assume that $T=\mathbb{Z}/p\mathbb{Z}\subseteq G$ is central, and apply it to finite $p$-groups $G$. \qed
\end{example}

\begin{example}
Assume that $G$ is a finite $p$-group, and $[G:T]=p$ (so $T$ is normal). \qed
\end{example}

